\documentclass[11pt,a4paper]{article}

\usepackage[utf8]{inputenc}
\usepackage[T1]{fontenc}
\usepackage[margin=1in]{geometry}
\usepackage{amsmath,amssymb,amsfonts,amsthm,mathtools}
\usepackage{bm}
\usepackage{booktabs}
\usepackage{array}
\usepackage{enumitem}
\usepackage{cite}
\usepackage[colorlinks=true,linkcolor=blue,citecolor=blue,urlcolor=blue]{hyperref}

% Theorem environments
\theoremstyle{plain}
\newtheorem{theorem}{Theorem}
\newtheorem{proposition}{Proposition}
\newtheorem{corollary}{Corollary}
\newtheorem{lemma}{Lemma}

\theoremstyle{definition}
\newtheorem{definition}{Definition}

\theoremstyle{remark}
\newtheorem{remark}{Remark}

% Custom mathematical macros
\newcommand{\Lcal}{\mathcal{L}}
\newcommand{\Kcal}{\mathcal{K}}
\newcommand{\Ical}{\mathcal{I}}
\newcommand{\Dcal}{\mathcal{D}}
\newcommand{\Rbb}{\mathbb{R}}
\newcommand{\dd}{\,\mathrm{d}}

\title{\textbf{Thermodynamic Consistency of Two-Phase Integral Nonlocal Thermoelasticity: Reciprocity, Boundary Closure, and Kernel Normalization}}

\author{Vipin Gupta}

\date{}

\begin{document}

\maketitle

\begin{abstract}
When Eringen's integral nonlocal thermoelasticity is converted into differential form via a Helmholtz kernel, care is required to preserve three properties of the underlying integral theory: (i)~Maxwell reciprocity between the thermal-stress operator and the dilational-entropy operator implied by the Second Law, (ii)~the Constitutive Boundary Conditions (CBCs) required for equivalence between the integral Fredholm operator and its differential inverse on bounded or semi-infinite domains, and (iii)~partition-of-unity kernel normalization near boundaries so that uniform temperature changes produce uniform free thermal expansion. Motivated by a convenience sample of $N=29$ nonlocal thermoelasticity papers published between 2016 and 2026 (drawn from a working collection of 152 mechanics PDFs, and including earlier papers co-authored by the present author), this clarification note consolidates the governing structure of two-phase ($\gamma \in (0,1)$) integral nonlocal thermoelasticity. First, specializing Eringen's (1974) integral kernel-duality theorem and the global nonlocality-residual framework to the two-phase operator $\Lcal_{\gamma,\ell}$, we state the two thermodynamically consistent realizations---Variant~V1 (thermo-mechanically nonlocal) and Variant~V2 (mechanically nonlocal)---and prove (\textbf{Theorem~1}) that the widely used differential shortcut pairing $(1-\ell^2\partial_x^2)\sigma = E\varepsilon - \beta\theta$ with the local constant-strain entropy $\rho\eta = \beta\varepsilon + (\rho c_E/T_0)\theta$ admits no twice-differentiable Helmholtz free energy in $(\varepsilon,\theta)$ and no stress-gradient Gibbs potential in $(\sigma,\theta)$ (where Gibbs reciprocity requires the additional stress-Laplacian term $(\beta\ell^2/E)\sigma_{,xx}$), producing negative internal dissipation on closed harmonic cycles. Second, we prove (\textbf{Theorem~2}) that Variant~V1 is eighth-order in space (quartic in $q=k^2$, $N=4$ wave-mode pairs) whereas Variant~V2 and the mismatched differential model are sixth-order (cubic in $q=k^2$, $N=3$ wave-mode pairs), and show that applying an identical single-phase ($\gamma=0$) nonlocal operator to stress, entropy, and heat flux cancels the nonlocal operator from the bulk heat equation via a singular perturbation limit in which the fourth root escapes as $q_4 \sim -1/(\gamma\ell^2)$. Third, extending isothermal nonlocal boundary closure results to thermoelasticity (\textbf{Theorem~3}), we derive the CBCs under unnormalized, Borino--Failla--Parrinello additive-defect normalized, and reflecting Neumann kernels, and show that a static $1/\sqrt{\gamma}$ surface thermal-expansion boundary layer under uniform heating is an artifact of unnormalized kernel truncation that vanishes identically ($\varepsilon^{\mathrm{V1}} \equiv \varepsilon^{\mathrm{V2}} \equiv \beta\theta_0/E$) under any partition-of-unity normalized symmetric kernel. Finally, an order-of-magnitude calibration for face-centered cubic Aluminum using published atomistic characteristic lengths ($\bar{\ell}_1 = 0.1203\;\mathrm{nm}$) shows that within the frequency band where Fourier heat conduction is physically defensible ($\omega\tau_q \le 0.1$), bulk thermoelastic damping differences between the formulations are $O((k\ell)^2) \lesssim 10^{-4}$. The value of this note is therefore mathematical and foundational rather than a prediction of a large measurable nanoscale bulk damping effect in monolithic crystals.
\end{abstract}

\noindent\textbf{Keywords:} Nonlocal thermoelasticity; Two-phase integral kernel; Clausius--Duhem inequality; Maxwell reciprocity; Constitutive boundary conditions; Kernel normalization.

\section{Introduction}\label{sec:intro}

When Eringen's integral nonlocal elasticity \cite{Eringen1972Linear,Eringen2002} is coupled with classical or generalized heat conduction to model nanoscale thermoelastic waves and resonators, the nonlocal attenuation operator is routinely replaced by the second-order Helmholtz differential operator $1 - \ell^2\partial_x^2$ (where $\ell = e_0 a$ is the internal characteristic length). For a homogeneous isotropic rod occupying a domain $\Omega \subseteq \Rbb$ with displacement $u(x,t)$, linearized strain $\varepsilon = \partial_x u$, and temperature increment $\theta = T - T_0$ ($|\theta|/T_0 \ll 1$), the balances of linear momentum and linearized entropy are
\begin{equation}\label{eq:balance_laws}
  \partial_x \sigma(x,t) + b(x,t) = \rho \ddot{u}(x,t), \qquad T_0 \rho \dot{\eta}(x,t) = -\partial_x q(x,t) + r(x,t),
\end{equation}
where $\sigma$ is the Cauchy stress, $\eta$ is the specific entropy density, $q$ is the heat flux, and $\rho, E, \beta, c_E, \kappa > 0$ denote the mass density, Young's modulus, thermal stress modulus, specific heat at constant strain, and thermal conductivity.

Although Eringen \cite{Eringen1974} derived the thermodynamically consistent integral equations of linear nonlocal thermoelasticity in 1974, subsequent applied boundary-value and wave-propagation studies have diverged into conflicting differential shortcuts regarding (i)~whether the nonlocal operator acts on the thermal stress $-\beta\theta$ without acting on the dilational heating $T_0\beta\dot{\varepsilon}$, (ii)~whether the purely nonlocal ($\gamma = 0$) differential operator is applied to all constitutive equations simultaneously, and (iii)~how the integral convolution is closed on bounded domains $\Omega = [0,L]$. This clarification note synthesizes the thermodynamic, spectral, and boundary-value conditions required to make two-phase nonlocal thermoelasticity well-posed and self-consistent.

\subsection{Prior Art and Theoretical Foundations}\label{subsec:prior_art}

Every mathematical ingredient used in this note rests on established prior art in nonlocal continuum mechanics and thermodynamics, which we credit explicitly:

\begin{enumerate}[leftmargin=*,align=left]
  \item \textbf{Edelen--Laws nonlocality residuals and Eringen's (1974) nonlocal thermoelasticity \cite{Eringen1972Linear,Eringen1974,Eringen2002}:} Building on the nonlocal thermodynamics of Edelen and Laws (1971), in which pointwise balance laws include zero-mean localization residuals ($\int_\Omega \hat{e}\dd x = 0$, $\int_\Omega \hat{h}\dd x = 0$), Eringen \cite{Eringen1974} (Section~3, Eqs.~(3.11)--(3.15)) proved that when the nonlocal Helmholtz free energy contains a bilinear thermomechanical coupling with a symmetric kernel $\alpha(|x-\xi|)$, the stress and entropy state laws must share the \emph{identical} integral operator:
  \begin{equation}\label{eq:eringen_1974_laws}
    t_{kl}(\bm{x},t) = \int_\Omega \Big[\lambda'(\lvert\bm{x}-\bm{\xi}\rvert)\varepsilon_{rr}(\bm{\xi},t)\delta_{kl} + 2\mu'(\lvert\bm{x}-\bm{\xi}\rvert)\varepsilon_{kl}(\bm{\xi},t) - \beta'(\lvert\bm{x}-\bm{\xi}\rvert)\theta(\bm{\xi},t)\Big]\dd V(\bm{\xi}),
  \end{equation}
  paired with $\rho\eta(\bm{x},t) = \int_\Omega \big[\beta'(\lvert\bm{x}-\bm{\xi}\rvert)\varepsilon_{kk}(\bm{\xi},t) + \frac{\rho c'}{T_0}(\lvert\bm{x}-\bm{\xi}\rvert)\theta(\bm{\xi},t)\big]\dd V(\bm{\xi})$. Thus, if the thermal stress is nonlocal with kernel $\beta'(|x-\xi|)$, the dilational entropy must carry the exact same kernel.

  \item \textbf{Polizzotto's variational thermodynamics and Borino--Failla--Parrinello (BFP) symmetric kernel normalization \cite{BazantJirasek2002,DiLuzio2005,MarottiDeSciarra2009,BazantLeHoover2010}:} Polizzotto (2001) reformulated Edelen--Laws nonlocal thermodynamics in Hilbert space via self-adjoint integral operators and variational principles. On a bounded domain $\Omega$, the unnormalized kernel $K_\ell^{\mathrm{unnorm}}(x,\xi) = \frac{1}{2\ell}\mathrm{e}^{-|x-\xi|/\ell}$ satisfies $\int_0^L K_\ell^{\mathrm{unnorm}}(x,\xi)\dd\xi = 1 - \vartheta_\ell(x) < 1$ with boundary weight defect $\vartheta_\ell(x) := \frac{1}{2}(\mathrm{e}^{-x/\ell} + \mathrm{e}^{-(L-x)/\ell})$, violating constant-field preservation near $\partial\Omega$, whereas naive row scaling $K_\ell(x,\xi)/\int_\Omega K_\ell(x,z)\dd z$ destroys symmetry and self-adjointness \cite{BazantJirasek2002,DiLuzio2005}. Borino, Failla, and Parrinello (2003) restored both symmetry and exact partition of unity ($\int_\Omega W_\ell^{\mathrm{BFP}}(x,\xi)\dd\xi \equiv 1$) by compensating the missing boundary weight with a local Dirac delta, $W_\ell^{\mathrm{BFP}}(x,\xi) = \vartheta_\ell(x)\delta(x-\xi) + K_\ell^{\mathrm{unnorm}}(x,\xi)$, as analyzed in detail by Di~Luzio \cite{DiLuzio2005}, Marotti de Sciarra \cite{MarottiDeSciarra2009}, and Ba{\v z}ant, Le, and Hoover \cite{BazantLeHoover2010}.

  \item \textbf{Challamel and Wang on boundary paradoxes of differential nonlocal beams \cite{ChallamelWang2014}:} Challamel and Wang (2008) and Challamel, Zhang, Wang, Reddy, Wang, Michelitsch, and Collet \cite{ChallamelWang2014} demonstrated that Eringen's differential relation $(1 - \ell^2\partial_x^2)M = -EIw''$ produces a non-self-adjoint operator under standard clamped-free boundary conditions ($\int_0^L (M''\delta w - w''\delta M)\dd x = [M'\delta M - M\delta M']_0^L/(EI\ell^{-2}) \neq 0$), and showed that self-adjointness and physical softening require either a discrete-continualized variational boundary condition or a two-phase local/nonlocal mixture.

  \item \textbf{Fern{\'a}ndez-S{\'a}ez, Zaera, Loya, and Reddy \cite{FernandezSaez2016,FernandezSaezZaera2017} and Romano, Barretta, Diaco, and Marotti de Sciarra \cite{Romano2017} on Constitutive Boundary Conditions (CBCs):} Fern{\'a}ndez-S{\'a}ez et al.\ \cite{FernandezSaez2016,FernandezSaezZaera2017} and Romano et al.\ \cite{Romano2017} proved that on a finite interval $[0,L]$, the integral relation $g(x) = \int_0^L K_\ell^{\mathrm{unnorm}}(x,\xi)f(\xi)\dd\xi$ is equivalent to the differential equation $(1-\ell^2\partial_x^2)g = f$ \emph{if and only if} $g$ additionally satisfies the two homogeneous Robin Constitutive Boundary Conditions $(\partial_x \mp \ell^{-1})g = 0$ at $x=0,L$. Without these CBCs, the purely nonlocal ($\gamma = 0$) integral problem is generally over-constrained and admits no smooth solution unless the source satisfies boundary compatibility conditions \cite{Romano2017}.

  \item \textbf{Mikhasev and Nobili \cite{MikhasevNobili2020} and Nobili and Pramanik \cite{NobiliPramanik2025a,NobiliPramanik2025b} on two-phase ($\gamma \in (0,1)$) singular perturbations:} Originally introduced in elasticity by Eringen (1987) and studied by Ba{\v z}ant and Jirásek \cite{BazantJirasek2002}, the two-phase operator $\Lcal_{\gamma,\ell}[f] = \gamma f + (1-\gamma)\Kcal_\ell[f]$ ($\gamma \in (0,1)$) regularizes the Fredholm integral equation of the first kind into one of the second kind. Mikhasev and Nobili \cite{MikhasevNobili2020} proved that as $\gamma \to 0^+$, the fourth-order differential reduction is a singular perturbation whose extra roots scale as $O(\gamma^{-1/2}\ell^{-1})$, spawning exponential boundary layers of thickness $\ell\sqrt{\gamma}$ near $x=0,L$. Nobili and Pramanik \cite{NobiliPramanik2025a,NobiliPramanik2025b} extended two-phase integral elasticity to viscoelasticity and demonstrated that boundary effects under unnormalized kernels mix true physical nonlocality with kernel truncation.

  \item \textbf{Giorgi and Morro \cite{GiorgiMorro2025} on thermodynamic consistency of nonlocal macroscopic models:} Giorgi and Morro \cite{GiorgiMorro2025} systematically analyzed thermodynamic consistency across nonlocal macroscopic theories using the global Clausius--Duhem inequality over $\Omega$, proving that cross-coupling terms between mechanical and thermal fields must arise as conjugate variational derivatives of a common thermodynamic potential.
\end{enumerate}

\subsection{Factual Survey of 29 Spatial-Nonlocal Thermoelasticity Papers}\label{subsec:survey}

To document how these foundational principles are applied in contemporary nonlocal thermoelasticity, we audited a convenience sample of $N = 29$ spatial-nonlocal thermoelasticity papers (published 2016--2026) drawn from a working collection of 152 mechanics PDFs. \textbf{Selection criterion:} Every paper in the 152-PDF collection that solves a spatial-nonlocal thermoelastic problem using Eringen's integral or differential nonlocal operator was included ($N = 29$; 25 solve coupled thermoelastic equations and 4 treat uncoupled thermal loads or uncoupled nonlocal heat conduction; 28 of the 29 are posed on bounded or half-space domains with boundaries). Because this collection was assembled around the author's research area, 13 of the 29 surveyed papers ($44.8\%$) are co-authored by the present author; it is therefore a convenience sample rather than a random sample of the literature.

Table~\ref{tab:census} classifies the 29 surveyed papers according to their constitutive pairing of the stress equation, entropy/energy equation, and boundary closure. All 25 coupled papers in the surveyed sample fall into two differential shortcut classes:

\begin{table}[htbp]
  \centering
  \caption{Classification of the $N = 29$ surveyed spatial-nonlocal thermoelasticity papers (2016--2026) by constitutive pairing and boundary treatment.}\label{tab:census}
  \vspace{4pt}
  \begin{tabular}{p{3.9cm}cccp{6.2cm}}
    \toprule
    \textbf{Constitutive Class} & \textbf{All ($N=29$)} & \textbf{Coupled ($N=25$)} & \textbf{Author Subset} & \textbf{Defining Equations \& Mathematical Status} \\
    \midrule
    \textbf{Class I:} Stress-nonlocal + local constant-strain heating & $12$ ($41.4\%$) & $12$ ($48.0\%$) & $2 / 12$ \cite{BarakGupta2023,GuptaKumar2023a} & $(1-\ell^2\nabla^2)\sigma = E\varepsilon - \beta\theta$ paired with $\rho\eta^{\mathrm{loc}} = \beta\varepsilon + \frac{\rho c_E}{T_0}\theta$. Violates Maxwell symmetry in both Helmholtz and Gibbs forms (Theorem~\ref{thm:no_free_energy}); omits Robin CBCs. \\
    \addlinespace[3pt]
    \textbf{Class II:} All-field identical-kernel cancellation ($\gamma=0$) & $13$ ($44.8\%$) & $13$ ($52.0\%$) & $11 / 13$ \cite{GuptaBarakDas2023b,GuptaBarakDas2024b} & Applies $(1-\ell^2\nabla^2)$ to $\sigma$, $\eta$, and $q$ identically. Cancels out of displacement--temperature PDEs in homogeneous domains; traction BCs become local if applied to $\sigma^{\mathrm{loc}}$. \\
    \addlinespace[3pt]
    \textbf{Class III:} Uncoupled thermal load / nonlocal DPL & $4$ ($13.8\%$) & $0$ ($0.0\%$) & $0 / 4$ & Prescribes $\theta(x,t)$ externally ($T_0\beta\dot{\varepsilon} = 0$) or applies $(1-\ell^2\nabla^2)$ only to uncoupled rigid-conductor heat conduction. \\
    \addlinespace[3pt]
    \textbf{Consistent two-phase (V1 or V2) with CBCs} & $0$ ($0.0\%$) & $0$ ($0.0\%$) & $0$ & Two-phase operator $\Lcal_{\gamma,\ell}$ ($\gamma\in(0,1)$) satisfying $\Lcal_\eta = \Lcal_{\mathrm{th}}^*$ with integral-consistent endpoint CBCs. \\
    \bottomrule
  \end{tabular}
\end{table}

\begin{itemize}[leftmargin=*]
  \item \textbf{Class I (Stress-nonlocal + local heating, 12 of 29 surveyed papers, $41.4\%$ of all, $48.0\%$ of coupled):} Applies $(1-\ell^2\nabla^2)\sigma = E\varepsilon - \beta\theta$ to the stress law while retaining the unmodified local entropy $\rho\eta^{\mathrm{loc}} = \beta\varepsilon + \frac{\rho c_E}{T_0}\theta$ in the energy equation. As shown in Theorem~\ref{thm:no_free_energy}, this hybrid pairing violates Eringen's \cite{Eringen1974} reciprocity identity \eqref{eq:eringen_1974_laws} and admits no free energy in either Helmholtz or Gibbs variables.
  \item \textbf{Class II (All-field identical-kernel cancellation, 13 of 29 surveyed papers, $44.8\%$ of all, $52.0\%$ of coupled):} Applies the identical purely nonlocal ($\gamma = 0$) differential operator $1 - \ell^2\nabla^2$ to the stress, entropy, and heat flux simultaneously. Substituting $(1-\ell^2\nabla^2)\sigma = E\varepsilon - \beta\theta$ into $\partial_x(1-\ell^2\nabla^2)\sigma = (1-\ell^2\nabla^2)(\rho\ddot{u})$ yields $(1-\ell^2\nabla^2)(E\partial_x^2 u - \beta\partial_x\theta - \rho\ddot{u}) = 0$, so the nonlocal factor cancels from the bulk displacement--temperature PDEs unless coupled through nonlocal Cauchy traction BCs and Robin CBCs.
\end{itemize}
We state factually that several of the present author's own earlier papers in the surveyed sample fall into these two classes (specifically, \cite{BarakGupta2023,GuptaKumar2023a} in Class~I and \cite{GuptaBarakDas2023b,GuptaBarakDas2024b} in Class~II). Clarifying how to formulate two-phase nonlocal thermoelasticity without falling into either Class~I or Class~II is the motivation for this note.

\subsection{Scope and Outline of This Clarification Note}\label{subsec:scope}

This note introduces no new physical postulates; it applies the established frameworks of Eringen \cite{Eringen1974}, Polizzotto and Borino--Failla--Parrinello \cite{DiLuzio2005,MarottiDeSciarra2009}, Fern{\'a}ndez-S{\'a}ez et al.\ \cite{FernandezSaez2016}, Romano et al.\ \cite{Romano2017}, Mikhasev and Nobili \cite{MikhasevNobili2020}, Nobili and Pramanik \cite{NobiliPramanik2025b}, and Giorgi and Morro \cite{GiorgiMorro2025} to resolve four specific questions:
\begin{enumerate}[leftmargin=*]
  \item \textbf{Section~\ref{sec:model}:} Formulates the general quadratic nonlocal Helmholtz free energy on $L^2(\Omega)$ and defines the two minimal thermodynamically consistent two-phase ($\gamma \in (0,1)$) realizations: \textbf{Variant~V1} (nonlocal elastic and thermal stress paired with nonlocal dilational entropy) and \textbf{Variant~V2} (nonlocal elastic stiffness paired with local thermal stress and local dilational entropy, as in Pham and Vu \cite{PhamVu2024}).
  \item \textbf{Section~\ref{sec:theorem1}:} Proves \textbf{Theorem~\ref{thm:no_free_energy}}, showing that the Class~I pairing admits no free energy in either Helmholtz $(\varepsilon,\theta)$ or Gibbs $(\sigma,\theta)$ variables, and provides a one-page hand-checkable Maxwell-symmetry and closed-cycle dissipation calculation.
  \item \textbf{Section~\ref{sec:mode_count}:} Proves \textbf{Theorem~\ref{thm:mode_count}}, showing that for $\gamma \in (0,1)$, Variant~V1 is an eighth-order spatial PDE system with $N=4$ wave-mode pairs $\pm k_j(\omega)$ whereas Variant~V2 and Class~I have $N=3$ pairs, and derives the singular-perturbation blow-up $k_4^2(\omega;\gamma) \sim -(\gamma\ell^2)^{-1}$ as $\gamma \to 0^+$.
  \item \textbf{Section~\ref{sec:boundary}:} Derives the exact endpoint Constitutive Boundary Conditions (\textbf{Theorem~\ref{thm:cbc}}) for the unnormalized exponential kernel $K_\ell^{\mathrm{unnorm}}$, the BFP-normalized kernel $W_\ell^{\mathrm{BFP}}$, and the Neumann Green's function $K_\ell^{\mathrm{Neu}}$, and proves (\textbf{Corollary~\ref{cor:artifact}}) that the static $1/\sqrt{\gamma}$ surface thermal-strain ratio between V2 and V1 under uniform heating is a kernel-truncation artifact of $K_\ell^{\mathrm{unnorm}}$ that vanishes identically under $W_\ell^{\mathrm{BFP}}$ and $K_\ell^{\mathrm{Neu}}$.
  \item \textbf{Section~\ref{sec:calibration}:} Evaluates the bulk attenuation differences for fcc Aluminum using the atomistic 2NN MEAM characteristic length $\bar{\ell}_1 = 0.1203\;\mathrm{nm}$ of Lazar, Agiasofitou, and B{\"o}hlke \cite{Lazar2022}, demonstrating plainly that within the frequency regime where Fourier heat conduction is physically valid, relative differences between V1, V2, and Class~I are $O(10^{-4})$ or smaller. Section~\ref{sec:limitations} summarizes the limitations and modeling guidelines, and Appendix~\ref{app:proofs} collects the full mathematical proofs.
\end{enumerate}

\section{Model and the Two Consistent Realizations V1 and V2}\label{sec:model}

\subsection{Two-Phase Integral Operator on $L^2(\Omega)$}\label{subsec:operators}

Let $\Omega \subseteq \Rbb$ be either the unbounded line $\Rbb$, the half-line $[0,\infty)$, or a bounded rod $[0,L]$ ($0 < L < \infty$). For a symmetric non-negative attenuation kernel $\Kcal_\ell : \Omega \times \Omega \to [0,\infty)$ satisfying $\Kcal_\ell(x,\xi) = \Kcal_\ell(\xi,x)$ with characteristic length $\ell > 0$, define the integral operator $\Kcal_\ell[f](x) := \int_\Omega \Kcal_\ell(x,\xi)f(\xi)\dd\xi$ and the \emph{two-phase local/nonlocal operator} \cite{BazantJirasek2002,MikhasevNobili2020,NobiliPramanik2025b} with local fraction $\gamma \in (0,1)$:
\begin{equation}\label{eq:L_gamma_def}
  \Lcal_{\gamma,\ell}[f](x) := \gamma f(x) + (1-\gamma)\int_\Omega \Kcal_\ell(x,\xi)f(\xi)\dd\xi = \big(\gamma\Ical + (1-\gamma)\Kcal_\ell\big)[f](x).
\end{equation}
Since $\Kcal_\ell(x,\xi) = \Kcal_\ell(\xi,x)$ and $\Kcal_\ell$ is positive-semidefinite on $L^2(\Omega)$, $\Lcal_{\gamma,\ell}$ is self-adjoint and strictly coercive on $L^2(\Omega)$ for every $\gamma \in (0,1)$:
\begin{equation}\label{eq:self_adjoint_coercive}
  \langle \Lcal_{\gamma,\ell}[f],\, g \rangle_{L^2(\Omega)} = \langle f,\, \Lcal_{\gamma,\ell}[g] \rangle_{L^2(\Omega)}, \qquad \langle \Lcal_{\gamma,\ell}[f],\, f \rangle_{L^2(\Omega)} \ge \gamma \lVert f \rVert_{L^2(\Omega)}^2 \qquad \forall\, f,g \in L^2(\Omega).
\end{equation}
On $\Omega = \Rbb$ with the bi-exponential Helmholtz kernel $K_\ell^{\mathrm{unnorm}}(x,\xi) := \frac{1}{2\ell}\mathrm{e}^{-|x-\xi|/\ell}$, where $(1-\ell^2\partial_x^2)K_\ell^{\mathrm{unnorm}}(x,\xi) = \delta(x-\xi)$, applying $1-\ell^2\partial_x^2$ to $g = \Lcal_{\gamma,\ell}[f]$ yields the differential form \cite{FernandezSaez2016,Romano2017,MikhasevNobili2020}:
\begin{equation}\label{eq:two_phase_diff}
  (1 - \ell^2\partial_x^2)g(x) = (1 - \gamma\ell^2\partial_x^2)f(x).
\end{equation}
On bounded domains $\Omega = [0,L]$, equivalence between the integral relation \eqref{eq:L_gamma_def} and the differential relation \eqref{eq:two_phase_diff} holds if and only if \eqref{eq:two_phase_diff} is supplemented by two endpoint Constitutive Boundary Conditions (derived for three kernel choices in Section~\ref{sec:boundary}).

\subsection{Nonlocal Clausius--Duhem Inequality and Operator Reciprocity}\label{subsec:reciprocity}

Following Eringen \cite{Eringen1974} (Section~2, Eqs.~(2.4)--(2.14)) and Polizzotto \cite{DiLuzio2005,MarottiDeSciarra2009}, let $\psi(x,t)$ be the specific Helmholtz free energy density and $\hat{e}(x,t)$, $\hat{h}(x,t)$ be the energy and entropy nonlocality residuals satisfying the global zero-production constraints $\int_\Omega \hat{e}\dd x = 0$ and $\int_\Omega \hat{h}\dd x = 0$. Integrating the pointwise dissipation density over $\Omega$ eliminates $\hat{e}$ and $\hat{h}$, yielding the global linearized Clausius--Duhem inequality \cite{Eringen1974,GiorgiMorro2025}:
\begin{equation}\label{eq:global_cd}
  \Dcal_{\mathrm{global}}(t) := -\frac{\dd}{\dd t}\Psi[\varepsilon,\theta] + \int_\Omega \left[\sigma(x,t)\dot{\varepsilon}(x,t) - \rho\eta(x,t)\dot{\theta}(x,t) - \frac{q(x,t)}{T_0}\partial_x\theta(x,t)\right]\dd x \ge 0,
\end{equation}
where $\Psi[\varepsilon,\theta] := \int_\Omega \rho\psi(x,t)\dd x$ is the total Helmholtz free energy functional. Consider the general linear nonlocal state laws on $L^2(\Omega)$:
\begin{equation}\label{eq:general_state_laws}
  \sigma = E\,\Lcal_{\mathrm{el}}[\varepsilon] - \beta\,\Lcal_{\mathrm{th}}[\theta], \qquad \rho\eta = \beta\,\Lcal_\eta[\varepsilon] + \frac{\rho c_E}{T_0}\,\Lcal_c[\theta],
\end{equation}
where $\Lcal_{\mathrm{el}}$, $\Lcal_{\mathrm{th}}$, $\Lcal_\eta$, and $\Lcal_c$ are bounded linear operators on $L^2(\Omega)$ with $\Lcal_{\mathrm{el}}^* = \Lcal_{\mathrm{el}}$ and $\Lcal_c^* = \Lcal_c$.

\begin{proposition}[Master Helmholtz Free Energy and Operator Reciprocity \cite{Eringen1974,GiorgiMorro2025}]\label{prop:reciprocity}
The linear state laws \eqref{eq:general_state_laws} are derivable as the Fr{\'e}chet derivatives
\begin{equation}\label{eq:frechet_state_laws}
  \sigma(x,t) = \frac{\delta\Psi[\varepsilon,\theta]}{\delta\varepsilon(x,t)}, \qquad \rho\eta(x,t) = -\frac{\delta\Psi[\varepsilon,\theta]}{\delta\theta(x,t)}
\end{equation}
of a twice Fr{\'e}chet-differentiable Helmholtz free energy functional $\Psi : L^2(\Omega)\times L^2(\Omega) \to \Rbb$ if and only if the thermal-stress and dilational-entropy operators are Hilbert-space adjoints of one another:
\begin{equation}\label{eq:reciprocity_condition}
  \Lcal_\eta = \Lcal_{\mathrm{th}}^* \qquad \text{on } L^2(\Omega).
\end{equation}
When $\Lcal_{\mathrm{th}}$ is self-adjoint ($\Lcal_{\mathrm{th}}^* = \Lcal_{\mathrm{th}}$), \eqref{eq:reciprocity_condition} reduces to $\Lcal_\eta = \Lcal_{\mathrm{th}}$, and the unique quadratic Helmholtz free energy functional (up to an additive constant) is
\begin{equation}\label{eq:master_psi}
  \Psi[\varepsilon,\theta] = \int_\Omega \left(\frac{1}{2}E\,\varepsilon\,\Lcal_{\mathrm{el}}[\varepsilon] - \beta\,\varepsilon\,\Lcal_{\mathrm{th}}[\theta] - \frac{\rho c_E}{2T_0}\,\theta\,\Lcal_c[\theta]\right)\dd x,
\end{equation}
for which the conservative mechanical--thermal power cancels identically in \eqref{eq:global_cd}, reducing $\Dcal_{\mathrm{global}}(t)$ to the purely conductive dissipation $\Dcal_{\mathrm{global}}(t) = -\int_\Omega \frac{q}{T_0}\partial_x\theta\dd x = \int_\Omega \frac{\kappa}{T_0}(\partial_x\theta)^2\dd x \ge 0$ under Fourier's law $q = -\kappa\partial_x\theta$.
\end{proposition}

\subsection{The Two Minimal Consistent Two-Phase Realizations: Variant V1 and Variant V2}\label{subsec:V1_V2}

In standard two-phase nonlocal thermoelasticity with elastic operator $\Lcal_{\mathrm{el}} = \Lcal_{\gamma,\ell}$, local thermal capacity $\Lcal_c = \Ical$, and Fourier heat conduction $q = -\kappa\partial_x\theta$ (with $r=0$), Proposition~\ref{prop:reciprocity} yields two distinct minimal consistent realizations depending on whether $\Lcal_{\mathrm{th}}$ is chosen equal to $\Lcal_{\gamma,\ell}$ or $\Ical$:

\paragraph{1. Variant V1 (Full Stress-Nonlocal and Entropy-Nonlocal Coupling: $\Lcal_{\mathrm{el}} = \Lcal_{\mathrm{th}} = \Lcal_\eta = \Lcal_{\gamma,\ell}$, $\Lcal_c = \Ical$).}
Applying the two-phase nonlocal operator $\Lcal_{\gamma,\ell}$ to the total local thermoelastic stress $\sigma^{\mathrm{loc}} = E\varepsilon - \beta\theta$ requires $\Lcal_{\mathrm{th}} = \Lcal_{\gamma,\ell}$, so reciprocity \eqref{eq:reciprocity_condition} mandates $\Lcal_\eta = \Lcal_{\gamma,\ell}$ \cite{Eringen1974}:
\begin{equation}\label{eq:V1_integral}
  \sigma^{\mathrm{V1}} = \Lcal_{\gamma,\ell}[E\varepsilon - \beta\theta], \qquad \rho\eta^{\mathrm{V1}} = \beta\,\Lcal_{\gamma,\ell}[\varepsilon] + \frac{\rho c_E}{T_0}\theta, \qquad q = -\kappa\partial_x\theta.
\end{equation}
Applying $(1-\ell^2\partial_x^2)$ via \eqref{eq:two_phase_diff} to both the stress law and the linearized entropy balance $\kappa\partial_x^2\theta = \rho c_E\dot{\theta} + T_0\beta\,\Lcal_{\gamma,\ell}[\dot{\varepsilon}]$ yields the coupled differential system of \textbf{Variant~V1}:
\begin{align}
  (1 - \ell^2\partial_x^2)\sigma^{\mathrm{V1}} &= (1 - \gamma\ell^2\partial_x^2)(E\partial_x u - \beta\theta), \qquad \partial_x\sigma^{\mathrm{V1}} + b = \rho\ddot{u}, \label{eq:V1_diff_stress}\\
  (1 - \ell^2\partial_x^2)(\kappa\partial_x^2\theta - \rho c_E\dot{\theta}) &= T_0\beta(1 - \gamma\ell^2\partial_x^2)\partial_x\dot{u}. \label{eq:V1_diff_heat}
\end{align}
Notice that $(1-\ell^2\partial_x^2)$ acts on $(\kappa\partial_x^2\theta - \rho c_E\dot{\theta})$ on the left-hand side of \eqref{eq:V1_diff_heat}, whereas $(1-\gamma\ell^2\partial_x^2)$ acts on $T_0\beta\partial_x\dot{u}$ on the right-hand side; for every $\gamma \in [0,1)$, these operators are distinct and do not cancel.

\paragraph{2. Variant V2 (Elastic-Only Nonlocal Coupling: $\Lcal_{\mathrm{el}} = \Lcal_{\gamma,\ell}$, $\Lcal_{\mathrm{th}} = \Lcal_\eta = \Lcal_c = \Ical$).}
If nonlocality is attributed strictly to interatomic elastic interactions while thermal pressure $-\beta\theta$ remains local ($\Lcal_{\mathrm{th}} = \Ical$), then reciprocity \eqref{eq:reciprocity_condition} mandates $\Lcal_\eta = \Ical$ \cite{PhamVu2024}:
\begin{equation}\label{eq:V2_integral}
  \sigma^{\mathrm{V2}} = \Lcal_{\gamma,\ell}[E\varepsilon] - \beta\theta, \qquad \rho\eta^{\mathrm{V2}} = \beta\varepsilon + \frac{\rho c_E}{T_0}\theta, \qquad q = -\kappa\partial_x\theta.
\end{equation}
Applying $(1-\ell^2\partial_x^2)$ to $\sigma^{\mathrm{V2}} + \beta\theta = \Lcal_{\gamma,\ell}[E\varepsilon]$ yields the coupled differential system of \textbf{Variant~V2}:
\begin{align}
  (1 - \ell^2\partial_x^2)\sigma^{\mathrm{V2}} &= E(1 - \gamma\ell^2\partial_x^2)\partial_x u - \beta(1 - \ell^2\partial_x^2)\theta, \qquad \partial_x\sigma^{\mathrm{V2}} + b = \rho\ddot{u}, \label{eq:V2_diff_stress}\\
  \kappa\partial_x^2\theta &= \rho c_E\dot{\theta} + T_0\beta\partial_x\dot{u}. \label{eq:V2_diff_heat}
\end{align}
Here the energy equation \eqref{eq:V2_diff_heat} retains its classical local form, while in \eqref{eq:V2_diff_stress} the operator $(1-\ell^2\partial_x^2)$ acts on $\beta\theta$ so that $\Lcal_{\gamma,\ell}^{-1}(1-\gamma\ell^2\partial_x^2)\theta = \theta$ remains local.

\section{Theorem 1: Incompatibility of Stress-Nonlocal + Local-Heating Pairing}\label{sec:theorem1}

In 12 of the 29 surveyed papers ($48.0\%$ of the 25 coupled thermoelastic papers), the stress equation of Variant~V1 ($\gamma=0$) is paired with the local entropy equation of Variant~V2:
\begin{equation}\label{eq:class_I_pairing}
  (1 - \ell^2\partial_x^2)\sigma = E\varepsilon - \beta\theta \quad \big(\text{or } \sigma = \Lcal_{\gamma,\ell}[E\varepsilon - \beta\theta]\big), \qquad \rho\eta^{\mathrm{loc}} = \beta\varepsilon + \frac{\rho c_E}{T_0}\theta.
\end{equation}
Theorem~\ref{thm:no_free_energy} proves that \eqref{eq:class_I_pairing} violates Maxwell reciprocity in both the strain--temperature (Helmholtz) and stress--temperature (Gibbs) representations.

\begin{theorem}[Non-Existence of Helmholtz or Gibbs Potentials for Class I]\label{thm:no_free_energy}
Let $\ell > 0$, $\gamma \in [0,1)$, $E > 0$, $\beta \neq 0$, and let $\Kcal_\ell \neq \Ical$ be any symmetric nonlocal kernel on $\Omega$. Consider the constitutive pairing
\begin{equation}\label{eq:thm1_stress}
  \sigma = \Lcal_{\gamma,\ell}[E\varepsilon - \beta\theta] \quad \Longleftrightarrow \quad (1 - \ell^2\partial_x^2)\sigma = (1 - \gamma\ell^2\partial_x^2)(E\varepsilon - \beta\theta),
\end{equation}
coupled with the local constant-strain entropy density and linearized energy balance
\begin{equation}\label{eq:thm1_entropy}
  \rho\eta^{\mathrm{loc}} = \beta\varepsilon + \frac{\rho c_E}{T_0}\theta \quad \Longrightarrow \quad -\partial_x q + r = T_0\rho\dot{\eta}^{\mathrm{loc}} = \rho c_E\dot{\theta} + T_0\beta\dot{\varepsilon}.
\end{equation}
Then:
\begin{enumerate}[label=(\alph*),leftmargin=*]
  \item \textbf{Non-existence of a Helmholtz free energy $\Psi[\varepsilon,\theta]$:} The mixed second Fr{\'e}chet derivatives of any candidate functional $\Psi[\varepsilon,\theta]$ on $L^2(\Omega)\times L^2(\Omega)$ have the non-vanishing Maxwell-symmetry defect operator
  \begin{equation}\label{eq:maxwell_defect_operator}
    \Delta_{\mathrm{Maxwell}}^{\mathrm{Helm}} := \frac{\delta\sigma}{\delta\theta} + \frac{\delta(\rho\eta^{\mathrm{loc}})}{\delta\varepsilon} = \beta\left(\Ical - \Lcal_{\gamma,\ell}\right) = \beta(1-\gamma)\left(\Ical - \Kcal_\ell\right) \neq 0 \quad \text{on } L^2(\Omega).
  \end{equation}
  Hence no twice Fr{\'e}chet-differentiable Helmholtz free energy $\Psi[\varepsilon,\theta]$ exists on $L^2(\Omega)\times L^2(\Omega)$.

  \item \textbf{Non-existence of a stress-gradient Gibbs potential $G[\sigma,\theta]$ ($\gamma=0$):} In the complementary stress-gradient $(\sigma,\theta)$ representation where $(1-\ell^2\partial_x^2)\sigma = E\varepsilon - \beta\theta$ is inverted as $\varepsilon = \frac{\delta G}{\delta\sigma} = \frac{1}{E}(\sigma - \ell^2\partial_x^2\sigma) + \frac{\beta}{E}\theta$ from the quadratic stress-gradient Gibbs functional
  \begin{equation}\label{eq:gibbs_functional}
    G[\sigma,\theta] = \int_\Omega \left[\frac{1}{2E}\sigma^2 + \frac{\ell^2}{2E}(\partial_x\sigma)^2 + \frac{\beta}{E}\sigma\theta + \frac{\rho c_\sigma}{2T_0}\theta^2\right]\dd x, \qquad \frac{\rho c_\sigma}{T_0} := \frac{\rho c_E}{T_0} + \frac{\beta^2}{E},
  \end{equation}
  the conjugate Gibbs entropy $\rho\eta^{\mathrm{Gibbs}} := \frac{\delta G}{\delta\theta}$ is uniquely determined to be
  \begin{equation}\label{eq:gibbs_entropy}
    \rho\eta^{\mathrm{Gibbs}} = \frac{\beta}{E}\sigma + \frac{\rho c_\sigma}{T_0}\theta = \beta\varepsilon + \frac{\rho c_E}{T_0}\theta + \frac{\beta\ell^2}{E}\partial_x^2\sigma = \rho\eta^{\mathrm{loc}} + \frac{\beta\ell^2}{E}\partial_x^2\sigma.
  \end{equation}
  Because $\rho\eta^{\mathrm{loc}}$ in \eqref{eq:thm1_entropy} omits the stress-Laplacian term $\frac{\beta\ell^2}{E}\partial_x^2\sigma \neq 0$, the pairing \eqref{eq:thm1_stress}--\eqref{eq:thm1_entropy} also violates Maxwell symmetry in the $(\sigma,\theta)$ Gibbs framework (compare the generalized thermodynamic potential analysis for higher-grade thermoelasticity by Forest and Amestoy \cite{ForestAmestoy2008}).

  \item \textbf{Negative internal dissipation on closed harmonic cycles:} Regardless of whether heat conduction is governed by Fourier, Cattaneo, Moore--Gibson--Thompson, or dual-phase-lag laws, there exist smooth time-periodic histories $(\varepsilon(x,t),\theta(x,t))$ of period $\mathcal{T} = 2\pi/\omega$ on $\Omega = [0,L]$ across which the net internal dissipation over a closed cycle is strictly negative: $\Delta\Dcal_{\mathrm{cycle}} := \oint \Dcal_{\mathrm{global}}(t)\dd t < 0$.
\end{enumerate}
\end{theorem}

\subsection{Explicit One-Page Hand-Check of Theorem \ref{thm:no_free_energy}}\label{subsec:hand_check}

Every claim of Theorem~\ref{thm:no_free_energy} can be verified directly by hand in three elementary steps on a single spatial Fourier mode $\cos(kx)$ ($k = n\pi/L > 0$, $n \in \Nbb$) on $\Omega = [0,L]$ with homogeneous Neumann conditions ($\partial_x\varepsilon = \partial_x\theta = \partial_x\sigma = 0$ at $x=0,L$), where $-\partial_x^2[\cos(kx)] = k^2\cos(kx)$:

\paragraph{Step 1: Modal Maxwell-symmetry defect in $(\varepsilon,\theta)$ variables.}
Set $\varepsilon(x,t) = \hat{\varepsilon}(t)\cos(kx)$ and $\theta(x,t) = \hat{\theta}(t)\cos(kx)$. Substituting into \eqref{eq:thm1_stress} and \eqref{eq:thm1_entropy} gives the algebraic modal relations
\begin{equation}\label{eq:hc_modal_laws}
  \hat{\sigma}(\hat{\varepsilon},\hat{\theta}) = \hat{\Lcal}_{\gamma,\ell}(k^2)\left(E\hat{\varepsilon} - \beta\hat{\theta}\right), \qquad \rho\hat{\eta}^{\mathrm{loc}}(\hat{\varepsilon},\hat{\theta}) = \beta\hat{\varepsilon} + \frac{\rho c_E}{T_0}\hat{\theta}, \qquad \hat{\Lcal}_{\gamma,\ell}(k^2) := \frac{1 + \gamma k^2\ell^2}{1 + k^2\ell^2}.
\end{equation}
If a twice-differentiable modal free energy $\hat{\psi}(\hat{\varepsilon},\hat{\theta})$ existed with $\hat{\sigma} = \partial\hat{\psi}/\partial\hat{\varepsilon}$ and $\rho\hat{\eta}^{\mathrm{loc}} = -\partial\hat{\psi}/\partial\hat{\theta}$, Schwarz's theorem ($\partial^2\hat{\psi}/\partial\hat{\theta}\partial\hat{\varepsilon} = \partial^2\hat{\psi}/\partial\hat{\varepsilon}\partial\hat{\theta}$) would require $\frac{\partial\hat{\sigma}}{\partial\hat{\theta}} + \frac{\partial(\rho\hat{\eta}^{\mathrm{loc}})}{\partial\hat{\varepsilon}} = 0$. Differentiating \eqref{eq:hc_modal_laws} gives
\begin{equation}\label{eq:hc_maxwell_defect}
  \frac{\partial\hat{\sigma}}{\partial\hat{\theta}} = -\beta\,\frac{1 + \gamma k^2\ell^2}{1 + k^2\ell^2}, \qquad \frac{\partial(\rho\hat{\eta}^{\mathrm{loc}})}{\partial\hat{\varepsilon}} = +\beta \quad\Longrightarrow\quad \frac{\partial\hat{\sigma}}{\partial\hat{\theta}} + \frac{\partial(\rho\hat{\eta}^{\mathrm{loc}})}{\partial\hat{\varepsilon}} = \frac{\beta(1-\gamma)k^2\ell^2}{1 + k^2\ell^2} \neq 0
\end{equation}
for all $k > 0$, $\ell > 0$, and $\gamma \in [0,1)$. By contrast, Variant~V1 ($\rho\hat{\eta}^{\mathrm{V1}} = \beta\hat{\Lcal}_{\gamma,\ell}(k^2)\hat{\varepsilon} + \frac{\rho c_E}{T_0}\hat{\theta}$) and Variant~V2 ($\hat{\sigma}^{\mathrm{V2}} = E\hat{\Lcal}_{\gamma,\ell}(k^2)\hat{\varepsilon} - \beta\hat{\theta}$, $\rho\hat{\eta}^{\mathrm{V2}} = \rho\hat{\eta}^{\mathrm{loc}}$) both satisfy $\frac{\partial\hat{\sigma}}{\partial\hat{\theta}} + \frac{\partial(\rho\hat{\eta})}{\partial\hat{\varepsilon}} = 0$ identically.

\paragraph{Step 2: Complementary Gibbs $(\sigma,\theta)$ check at $\gamma = 0$.}
From \eqref{eq:gibbs_functional} with integrand $g(\sigma,\partial_x\sigma,\theta)$, the Euler--Lagrange derivative with respect to $\sigma$ (with $\partial_x\sigma = 0$ at $x=0,L$) gives $\varepsilon = \frac{\partial g}{\partial\sigma} - \partial_x\frac{\partial g}{\partial(\partial_x\sigma)} = \frac{1}{E}(\sigma - \ell^2\partial_x^2\sigma) + \frac{\beta}{E}\theta$, which rearranges to $(1-\ell^2\partial_x^2)\sigma = E\varepsilon - \beta\theta$, i.e., $\sigma = E\varepsilon - \beta\theta + \ell^2\partial_x^2\sigma$. Substituting this expression for $\sigma$ into $\rho\eta^{\mathrm{Gibbs}} = \frac{\partial g}{\partial\theta} = \frac{\beta}{E}\sigma + \frac{\rho c_\sigma}{T_0}\theta$ gives
\begin{equation}\label{eq:hc_gibbs}
  \rho\eta^{\mathrm{Gibbs}} = \frac{\beta}{E}\left(E\varepsilon - \beta\theta + \ell^2\partial_x^2\sigma\right) + \left(\frac{\rho c_E}{T_0} + \frac{\beta^2}{E}\right)\theta = \underbrace{\beta\varepsilon + \frac{\rho c_E}{T_0}\theta}_{\rho\eta^{\mathrm{loc}}} + \frac{\beta\ell^2}{E}\partial_x^2\sigma.
\end{equation}
The terms $-\frac{\beta^2}{E}\theta$ and $+\frac{\beta^2}{E}\theta$ cancel identically, leaving the non-zero discrepancy $\rho\eta^{\mathrm{Gibbs}} - \rho\eta^{\mathrm{loc}} = \frac{\beta\ell^2}{E}\partial_x^2\sigma \neq 0$.

\paragraph{Step 3: Closed harmonic cycle yielding negative internal dissipation $\Delta\Dcal_{\mathrm{cycle}} < 0$.}
Prescribe the periodic history $\varepsilon(x,t) = \varepsilon_0\cos(kx)\cos(\omega t)$, $\theta(x,t) = \theta_0\cos(kx)\sin(\omega t)$ ($\varepsilon_0,\theta_0,\omega > 0$) over one period $\mathcal{T} = 2\pi/\omega$, where $\oint \dot{\Psi}\dd t = 0$ for any state function $\Psi$. Using $\int_0^L \cos^2(kx)\dd x = \frac{L}{2}$, $\int_0^{\mathcal{T}}\cos(\omega t)\sin(\omega t)\dd t = 0$, and $\int_0^{\mathcal{T}}\sin^2(\omega t)\dd t = \int_0^{\mathcal{T}}\cos^2(\omega t)\dd t = \frac{\pi}{\omega}$:
\begin{align}
  \oint_{0}^{\mathcal{T}}\!\!\int_0^L \sigma\dot{\varepsilon}\dd x\dd t &= \frac{L}{2}\int_0^{\mathcal{T}} \hat{\Lcal}_{\gamma,\ell}(k^2)\big(E\varepsilon_0\cos\omega t - \beta\theta_0\sin\omega t\big)(-\omega\varepsilon_0\sin\omega t)\dd t = +\frac{\pi L}{2}\beta\varepsilon_0\theta_0\,\hat{\Lcal}_{\gamma,\ell}(k^2), \label{eq:hc_cycle_work}\\
  -\oint_{0}^{\mathcal{T}}\!\!\int_0^L \rho\eta^{\mathrm{loc}}\dot{\theta}\dd x\dd t &= -\frac{L}{2}\int_0^{\mathcal{T}} \Big(\beta\varepsilon_0\cos\omega t + \frac{\rho c_E}{T_0}\theta_0\sin\omega t\Big)(\omega\theta_0\cos\omega t)\dd t = -\frac{\pi L}{2}\beta\varepsilon_0\theta_0. \label{eq:hc_cycle_entropy}
\end{align}
Adding the Fourier thermal conduction integral $\oint_0^{\mathcal{T}}\int_0^L \frac{\kappa}{T_0}(\partial_x\theta)^2\dd x\dd t = \frac{\pi L}{2\omega}\frac{\kappa k^2}{T_0}\theta_0^2$ to \eqref{eq:hc_cycle_work}--\eqref{eq:hc_cycle_entropy} yields
\begin{equation}\label{eq:hc_cycle_total}
  \Delta\Dcal_{\mathrm{cycle}} = \frac{\pi L\theta_0}{2\omega}\left[-\frac{\beta(1-\gamma)k^2\ell^2}{1 + k^2\ell^2}\,\omega\varepsilon_0 + \frac{\kappa k^2}{T_0}\,\theta_0\right] < 0 \quad \text{whenever}\quad \varepsilon_0 > \frac{\kappa(1 + k^2\ell^2)}{\beta T_0(1-\gamma)\omega\ell^2}\,\theta_0.
\end{equation}
Thus the mismatched pairing \eqref{eq:class_I_pairing} acts as an unphysical internal energy source during high-frequency thermomechanical cycles, whereas for both Variant~V1 and Variant~V2 the first bracketed term in \eqref{eq:hc_cycle_total} vanishes identically, leaving $\Delta\Dcal_{\mathrm{cycle}}^{\mathrm{V1}} = \Delta\Dcal_{\mathrm{cycle}}^{\mathrm{V2}} = \frac{\pi L}{2\omega}\frac{\kappa k^2}{T_0}\theta_0^2 > 0$.

\section{Mode-Count Theorem ($N=3$ vs.\ $N=4$) and the $\gamma \to 0$ Singular Limit}\label{sec:mode_count}

We now show that enforcing the nonlocal dilational entropy law in Variant~V1 changes the spatial differential order of the coupled thermoelastic system from sixth order ($N=3$ mode pairs $\pm k_j(\omega)$) to eighth order ($N=4$ mode pairs), and that the purely nonlocal limit $\gamma \to 0^+$ is a singular perturbation \cite{MikhasevNobili2020}.

\subsection{Master Dispersion Equation and Spatial Polynomial Degree}\label{subsec:master_dispersion}

Seek harmonic plane-wave solutions $(u,\theta)(x,t) = (\hat{u},\hat{\theta})\mathrm{e}^{\mathrm{i}(kx-\omega t)}$ ($\omega > 0$) on $\Omega = \Rbb$ with $q := k^2 \in \Cbb$ and Fourier symbol
\begin{equation}\label{eq:L_symbol}
  \hat{\Lcal}_{\gamma,\ell}(q) = \gamma + \frac{1-\gamma}{1+\ell^2 q} = \frac{1 + \gamma\ell^2 q}{1 + \ell^2 q}.
\end{equation}
For the general state laws \eqref{eq:general_state_laws} with $\Lcal_c = \Ical$ and Fourier conduction, the $2\times 2$ homogeneous system for $(\hat{u},\hat{\theta})$ is
\begin{equation}\label{eq:harmonic_matrix}
  \begin{bmatrix}
    E q\,\hat{\Lcal}_{\mathrm{el}}(q) - \rho\omega^2 & \mathrm{i}k\beta\,\hat{\Lcal}_{\mathrm{th}}(q) \\[4pt]
    -\omega k T_0\beta\,\hat{\Lcal}_\eta(q) & \kappa q - \mathrm{i}\omega\rho c_E
  \end{bmatrix}
  \begin{bmatrix}
    \hat{u} \\ \hat{\theta}
  \end{bmatrix}
  = \begin{bmatrix} 0 \\ 0 \end{bmatrix},
\end{equation}
yielding the master rational dispersion equation
\begin{equation}\label{eq:master_dispersion}
  \left(E q\,\hat{\Lcal}_{\mathrm{el}}(q) - \rho\omega^2\right)\left(\kappa q - \mathrm{i}\omega\rho c_E\right) - \mathrm{i}\omega T_0\beta^2 q\,\hat{\Lcal}_{\mathrm{th}}(q)\hat{\Lcal}_\eta(q) = 0.
\end{equation}

\begin{theorem}[Mode Count $N=3$ vs.\ $N=4$ and Singular Limit $\gamma \to 0^+$]\label{thm:mode_count}
Let $E,\rho,c_E,\kappa,T_0,\beta,\ell > 0$, $\gamma \in (0,1)$, and $\omega > 0$. Clearing the denominator $(1+\ell^2 q)$ in \eqref{eq:master_dispersion} yields the following polynomial dispersion equations $P(q,\omega) = 0$ in $q = k^2$:
\begin{enumerate}[label=(\alph*),leftmargin=*]
  \item \textbf{Exact spatial mode counts:}
  \begin{itemize}[leftmargin=*]
    \item \textbf{Variant V1} ($\hat{\Lcal}_{\mathrm{el}} = \hat{\Lcal}_{\mathrm{th}} = \hat{\Lcal}_\eta = \hat{\Lcal}_{\gamma,\ell}$): Because $\hat{\Lcal}_{\mathrm{th}}\hat{\Lcal}_\eta = (1+\gamma\ell^2 q)^2/(1+\ell^2 q)^2$ carries a squared denominator $(1+\ell^2 q)^2$, multiplying \eqref{eq:master_dispersion} by $(1+\ell^2 q)^2$ yields a \textbf{quartic polynomial} in $q = k^2$ ($N=4$ wave-mode pairs $\pm k_1,\pm k_2,\pm k_3,\pm k_4$, eighth-order PDE in $x$):
    \begin{equation}\label{eq:P_V1}
      P_{\mathrm{V1}}(q,\omega;\gamma) := \big[E q(1+\gamma\ell^2 q) - \rho\omega^2(1+\ell^2 q)\big](\kappa q - \mathrm{i}\omega\rho c_E)(1+\ell^2 q) - \mathrm{i}\omega T_0\beta^2 q(1+\gamma\ell^2 q)^2 = 0,
    \end{equation}
    with non-vanishing leading coefficient $a_4^{\mathrm{V1}} = \gamma\ell^4 E\kappa > 0$ and no root at $q = -\ell^{-2}$ ($P_{\mathrm{V1}}(-\ell^{-2},\omega;\gamma) = \mathrm{i}\omega T_0\beta^2\ell^{-2}(1-\gamma)^2 \neq 0$).
    \item \textbf{Variant V2} ($\hat{\Lcal}_{\mathrm{el}} = \hat{\Lcal}_{\gamma,\ell}$, $\hat{\Lcal}_{\mathrm{th}} = \hat{\Lcal}_\eta = 1$): Multiplying \eqref{eq:master_dispersion} by $(1+\ell^2 q)$ yields a \textbf{cubic polynomial} in $q = k^2$ ($N=3$ wave-mode pairs $\pm k_1,\pm k_2,\pm k_3$, sixth-order PDE in $x$):
    \begin{equation}\label{eq:P_V2}
      P_{\mathrm{V2}}(q,\omega;\gamma) := \big[E q(1+\gamma\ell^2 q) - \rho\omega^2(1+\ell^2 q)\big](\kappa q - \mathrm{i}\omega\rho c_E) - \mathrm{i}\omega T_0\beta^2 q(1+\ell^2 q) = 0,
    \end{equation}
    with leading coefficient $a_3^{\mathrm{V2}} = \gamma\ell^2 E\kappa > 0$.
    \item \textbf{Class I} ($\gamma = 0$ stress-nonlocal + local heating: $\hat{\Lcal}_{\mathrm{el}} = \hat{\Lcal}_{\mathrm{th}} = (1+\ell^2 q)^{-1}$, $\hat{\Lcal}_\eta = 1$): Multiplying \eqref{eq:master_dispersion} by $(1+\ell^2 q)$ yields a \textbf{cubic polynomial} in $q = k^2$ ($N=3$ wave-mode pairs):
    \begin{equation}\label{eq:P_ClassI}
      P_{\mathrm{Class\,I}}(q,\omega) := \big[E q - \rho\omega^2(1+\ell^2 q)\big](\kappa q - \mathrm{i}\omega\rho c_E) - \mathrm{i}\omega T_0\beta^2 q = 0,
    \end{equation}
    whose leading coefficient $a_3^{\mathrm{Class\,I}}(\omega) = \kappa(E - \rho\omega^2\ell^2)$ changes sign at the cut-off frequency $\omega_c = c_0/\ell$ ($c_0 := \sqrt{E/\rho}$).
  \end{itemize}

  \item \textbf{Singular limit $\gamma \to 0^+$ of Variant V1:} As $\gamma \to 0^+$ (for $\omega \neq c_0/\ell$), three roots $q_1(\omega;\gamma), q_2(\omega;\gamma), q_3(\omega;\gamma)$ of \eqref{eq:P_V1} converge smoothly to the three roots of the $\gamma = 0$ Eringen cubic
  \begin{equation}\label{eq:P_V1_gamma0}
    P_{\mathrm{V1}}(q,\omega;0) = \big[E q - \rho\omega^2(1+\ell^2 q)\big](\kappa q - \mathrm{i}\omega\rho c_E)(1+\ell^2 q) - \mathrm{i}\omega T_0\beta^2 q = 0,
  \end{equation}
  while the fourth root $q_4(\omega;\gamma)$ diverges to infinity as $O(\gamma^{-1}\ell^{-2})$ with the asymptotic expansion
  \begin{equation}\label{eq:q4_asymptotic}
    q_4(\omega;\gamma) = -\frac{1}{\gamma\ell^2} + \frac{\rho\omega^2}{E} - \frac{\mathrm{i}\omega T_0\beta^2}{\kappa E}\,\gamma + O(\gamma^2) \qquad (\gamma \to 0^+),
  \end{equation}
  generating an evanescent boundary-layer pair $k_4(\omega;\gamma) = \pm\frac{\mathrm{i}}{\ell\sqrt{\gamma}}\big[1 - \frac{\gamma\rho\omega^2\ell^2}{2E} + O(\gamma^2)\big]$ of decay length $\delta_{\mathrm{BL}}(\gamma) = \ell\sqrt{\gamma} \to 0$.

  \item \textbf{Asymptotic bulk thermoelastic attenuation ($\alpha_1 := \mathrm{Im}\,k_1(\omega)$):} Define the dimensionless thermoelastic coupling $\epsilon := \frac{T_0\beta^2}{\rho c_E E}$, characteristic frequency $\omega^* := \frac{\rho c_E c_0^2}{\kappa}$, characteristic length $\ell^* := \frac{c_0}{\omega^*}$, and dimensionless frequency $\Omega := \omega/\omega^*$.
  \begin{enumerate}[label=(\roman*)]
    \item \emph{Long-wavelength regime ($k\ell \ll 1$, $\epsilon \ll 1$):} All three models agree to leading local order $\alpha_1^{\mathrm{loc}}(\Omega) = \frac{\epsilon}{2\ell^*}\frac{\Omega^2}{1+\Omega^2}$ and differ at $O((k\ell)^2) = O((\ell/\ell^*)^2\Omega^2)$ with distinct numerical prefactors:
    \begin{equation}\label{eq:low_k_damping_expansion}
      \frac{\alpha_1(\Omega)}{\alpha_1^{\mathrm{loc}}(\Omega)} = 1 + (1-\gamma)(k_0\ell)^2 \times \begin{cases}
        -\dfrac{1}{2}, & \text{Variant V1}, \\[6pt]
        +\dfrac{3}{2}, & \text{Variant V2}, \\[6pt]
        +\dfrac{1}{2}\; (\text{at }\gamma=0), & \text{Class I},
      \end{cases}
      \quad + O\big(\epsilon (k_0\ell)^2, (k_0\ell)^4\big),
    \end{equation}
    where $k_0 := \omega/c_0$.
    \item \emph{Short-wavelength limit ($\Omega \to \infty$, $\gamma \in (0,1)$):} The high-frequency quasi-elastic attenuation plateaus of Variant~V1 and Variant~V2 differ by the exact factor $\gamma^2$:
    \begin{equation}\label{eq:high_freq_plateau}
      \lim_{\Omega\to\infty} \ell^*\alpha_1^{\mathrm{V1}}(\Omega) = \frac{\epsilon}{2}\gamma^{3/2} + O(\epsilon^2), \qquad
      \lim_{\Omega\to\infty} \ell^*\alpha_1^{\mathrm{V2}}(\Omega) = \frac{\epsilon}{2}\gamma^{-1/2} + O(\epsilon^2), \qquad
      \lim_{\Omega\to\infty}\frac{\alpha_1^{\mathrm{V1}}(\Omega)}{\alpha_1^{\mathrm{V2}}(\Omega)} = \gamma^2 + O(\epsilon).
    \end{equation}
  \end{enumerate}
\end{enumerate}
\end{theorem}

\section{Boundary Closure Under Unnormalized, BFP-Normalized, and Neumann Kernels}\label{sec:boundary}

On a bounded rod $\Omega = [0,L]$, converting the two-phase integral relation $g(x) = \Lcal_{\gamma,\ell}[f](x)$ into a differential equation requires specifying the exact kernel $\Kcal_\ell(x,\xi)$ on $[0,L]^2$ and its endpoint Constitutive Boundary Conditions \cite{FernandezSaez2016,Romano2017,MikhasevNobili2020,NobiliPramanik2025b}.

\subsection{Three Symmetric Kernels on $[0,L]$: Symmetry vs.\ Partition of Unity}\label{subsec:three_kernels}

We compare three symmetric ($\Kcal_\ell(x,\xi) = \Kcal_\ell(\xi,x)$), positive-semidefinite kernels on $[0,L]^2$:

\begin{enumerate}[leftmargin=*,align=left]
  \item \textbf{Unnormalized bi-exponential kernel ($K_\ell^{\mathrm{unnorm}}$) \cite{FernandezSaez2016,Romano2017,MikhasevNobili2020}:}
  \begin{equation}\label{eq:K_unnorm}
    K_\ell^{\mathrm{unnorm}}(x,\xi) := \frac{1}{2\ell}\mathrm{e}^{-|x-\xi|/\ell}, \qquad
    \int_0^L K_\ell^{\mathrm{unnorm}}(x,\xi)\dd\xi = 1 - \vartheta_\ell(x) < 1,
  \end{equation}
  where $\vartheta_\ell(x) := \frac{1}{2}\big(\mathrm{e}^{-x/\ell} + \mathrm{e}^{-(L-x)/\ell}\big) > 0$ is the boundary weight defect lost to $\Rbb \setminus [0,L]$. Because $\vartheta_\ell(x) > 0$, $K_\ell^{\mathrm{unnorm}}$ does not preserve constant fields ($\Lcal_{\gamma,\ell}^{\mathrm{unnorm}}[1](x) = 1 - (1-\gamma)\vartheta_\ell(x) \neq 1$).

  \item \textbf{Borino--Failla--Parrinello (BFP) additive-defect symmetric normalized kernel ($W_\ell^{\mathrm{BFP}}$) \cite{DiLuzio2005,MarottiDeSciarra2009,BazantLeHoover2010}:}
  \begin{equation}\label{eq:W_BFP}
    W_\ell^{\mathrm{BFP}}(x,\xi) := \vartheta_\ell(x)\delta(x-\xi) + K_\ell^{\mathrm{unnorm}}(x,\xi), \qquad
    \int_0^L W_\ell^{\mathrm{BFP}}(x,\xi)\dd\xi \equiv 1.
  \end{equation}
  The corresponding two-phase operator has a spatially varying effective local fraction $\gamma_{\mathrm{eff}}(x) := \gamma + (1-\gamma)\vartheta_\ell(x)$:
  \begin{equation}\label{eq:L_BFP}
    \Lcal_{\gamma,\ell}^{\mathrm{BFP}}[f](x) = \gamma_{\mathrm{eff}}(x)f(x) + (1-\gamma)\int_0^L K_\ell^{\mathrm{unnorm}}(x,\xi)f(\xi)\dd\xi, \qquad
    \Lcal_{\gamma,\ell}^{\mathrm{BFP}}[1](x) \equiv 1.
  \end{equation}

  \item \textbf{Neumann (cos-reflective) Green's function kernel ($K_\ell^{\mathrm{Neu}}$) \cite{MarottiDeSciarra2009,BazantLeHoover2010}:} Let $K_\ell^{\mathrm{Neu}}(x,\xi)$ be the Green's function of $(1-\ell^2\partial_x^2)$ on $[0,L]$ with homogeneous Neumann boundary conditions $\partial_x K_\ell^{\mathrm{Neu}}(0,\xi) = \partial_x K_\ell^{\mathrm{Neu}}(L,\xi) = 0$:
  \begin{equation}\label{eq:K_Neu}
    K_\ell^{\mathrm{Neu}}(x,\xi) := \frac{\cosh\!\big((L - |x-\xi|)/\ell\big) + \cosh\!\big((L - (x+\xi))/\ell\big)}{2\ell\sinh(L/\ell)}, \qquad
    \int_0^L K_\ell^{\mathrm{Neu}}(x,\xi)\dd\xi \equiv 1,
  \end{equation}
  so that $\Lcal_{\gamma,\ell}^{\mathrm{Neu}}[1](x) \equiv 1$ and $K_\ell^{\mathrm{Neu}}(x,\xi) = K_\ell^{\mathrm{unnorm}}(x,\xi) + \frac{1}{2\ell}\mathrm{e}^{-(x+\xi)/\ell} + O(\mathrm{e}^{-L/\ell})$ via specular reflection across $x=0,L$.
\end{enumerate}

\subsection{Theorem 3: Differential Equivalence and Endpoint CBCs}\label{subsec:thm3_cbc}

\begin{theorem}[Exact Differential Equivalence and Constitutive Boundary Conditions on ${[0,L]}$]\label{thm:cbc}
Let $\gamma \in (0,1)$, $\ell > 0$, and $f \in H^2(0,L)$. A function $g \in H^2(0,L)$ satisfies the two-phase integral relation $g = \Lcal_{\gamma,\ell}[f]$ on $[0,L]$ if and only if $g$ satisfies the following interior differential equation on $(0,L)$ and two endpoint Constitutive Boundary Conditions (CBCs) at $x \in \{0,L\}$:
\begin{enumerate}[label=(\alph*),leftmargin=*]
  \item \textbf{Under the unnormalized kernel $K_\ell^{\mathrm{unnorm}}$ \eqref{eq:K_unnorm} \cite{FernandezSaez2016,Romano2017,MikhasevNobili2020}:}
  \begin{align}
    (1 - \ell^2\partial_x^2)g(x) &= (1 - \gamma\ell^2\partial_x^2)f(x), && x \in (0,L), \label{eq:ode_unnorm}\\
    \left(\partial_x - \frac{1}{\ell}\right)\big(g - \gamma f\big)\Big|_{x=0} &= 0, \qquad
    \left(\partial_x + \frac{1}{\ell}\right)\big(g - \gamma f\big)\Big|_{x=L} = 0. \label{eq:cbc_unnorm}
  \end{align}

  \item \textbf{Under the BFP-normalized kernel $W_\ell^{\mathrm{BFP}}$ \eqref{eq:W_BFP}:} With $\gamma_{\mathrm{eff}}(x) = \gamma + (1-\gamma)\vartheta_\ell(x)$,
  \begin{align}
    (1 - \ell^2\partial_x^2)\big(g(x) - \gamma_{\mathrm{eff}}(x)f(x)\big) &= (1 - \gamma)f(x), && x \in (0,L), \label{eq:ode_bfp}\\
    \left(\partial_x - \frac{1}{\ell}\right)\big(g - \gamma_{\mathrm{eff}}f\big)\Big|_{x=0} &= 0, \qquad
    \left(\partial_x + \frac{1}{\ell}\right)\big(g - \gamma_{\mathrm{eff}}f\big)\Big|_{x=L} = 0. \label{eq:cbc_bfp}
  \end{align}

  \item \textbf{Under the Neumann Green's function kernel $K_\ell^{\mathrm{Neu}}$ \eqref{eq:K_Neu}:}
  \begin{align}
    (1 - \ell^2\partial_x^2)g(x) &= (1 - \gamma\ell^2\partial_x^2)f(x), && x \in (0,L), \label{eq:ode_neu}\\
    \partial_x\big(g - \gamma f\big)\Big|_{x=0} &= 0, \qquad
    \partial_x\big(g - \gamma f\big)\Big|_{x=L} = 0. \label{eq:cbc_neu}
  \end{align}
\end{enumerate}
\end{theorem}

\paragraph{Complete boundary closure for Variants V1 and V2.}
On $\Omega = [0,L]$, a well-posed boundary-value problem combines four classical physical boundary conditions at $x \in \{0,L\}$ (two mechanical conditions on $u$ or the nonlocal Cauchy traction $\sigma$, never on $\sigma^{\mathrm{loc}} = E\varepsilon - \beta\theta$, plus two thermal conditions on $\theta$ or $q = -\kappa\partial_x\theta$) with the CBCs of Theorem~\ref{thm:cbc}:
\begin{itemize}[leftmargin=*]
  \item \textbf{Variant V2 ($N=3$, sixth-order in $x$):} Requires $2N = 6$ boundary conditions, provided by the 4 physical BCs together with the 2 stress CBCs of Theorem~\ref{thm:cbc} evaluated for $g = \sigma^{\mathrm{V2}} + \beta\theta$ and $f = E\varepsilon$.
  \item \textbf{Variant V1 ($N=4$, eighth-order in $x$):} Requires $2N = 8$ boundary conditions, provided by the 4 physical BCs together with 4 CBCs: the 2 stress CBCs ($g = \sigma^{\mathrm{V1}}$, $f = E\varepsilon - \beta\theta$) and the 2 entropy CBCs ($g = \rho\eta^{\mathrm{V1}} - \frac{\rho c_E}{T_0}\theta$, $f = \beta\varepsilon$).
\end{itemize}

\subsection{The Static $1/\sqrt{\gamma}$ Surface Thermal-Strain Spike Is a Truncation Artifact}\label{subsec:truncation_artifact}

Consider a traction-free rod ($\sigma(0) = \sigma(L) = 0$, $b \equiv 0 \implies \sigma(x) \equiv 0$ on $[0,L]$) held at a uniform static temperature rise $\theta(x) \equiv \theta_0 \neq 0$. Under Variant~V1 ($\sigma^{\mathrm{V1}} = \Lcal_{\gamma,\ell}[E\varepsilon^{\mathrm{V1}} - \beta\theta_0] = 0$), coercivity of $\Lcal_{\gamma,\ell}$ on $L^2(0,L)$ (Eq.~\eqref{eq:self_adjoint_coercive}) uniquely forces $E\varepsilon^{\mathrm{V1}}(x) - \beta\theta_0 \equiv 0$, so $\varepsilon^{\mathrm{V1}}(x) \equiv \frac{\beta\theta_0}{E}$ is spatially uniform for \emph{any} coercive kernel.

Under Variant~V2 ($\sigma^{\mathrm{V2}} = \Lcal_{\gamma,\ell}[E\varepsilon^{\mathrm{V2}}] - \beta\theta_0 = 0$), the strain solves $\Lcal_{\gamma,\ell}[E\varepsilon^{\mathrm{V2}}](x) = \beta\theta_0$. When the unnormalized kernel $K_\ell^{\mathrm{unnorm}}$ is used, Theorem~\ref{thm:cbc}(a) with $g(x) \equiv \beta\theta_0$ and $f(x) = E\varepsilon^{\mathrm{V2}}(x)$ yields $(1-\gamma\ell^2\partial_x^2)(E\varepsilon^{\mathrm{V2}}) = \beta\theta_0$ with Robin CBCs $(\partial_x \mp \ell^{-1})(\beta\theta_0 - \gamma E\varepsilon^{\mathrm{V2}})\big|_{x=0,L} = 0$, whose exact solution on $[0,L]$ is
\begin{equation}\label{eq:unnorm_uniform_strain}
  \varepsilon^{\mathrm{V2,unnorm}}(x) = \frac{\beta\theta_0}{E}\left[1 + \frac{1-\gamma}{\sqrt{\gamma}}\,\frac{\cosh\!\left(\dfrac{x - L/2}{\ell\sqrt{\gamma}}\right)}{\sqrt{\gamma}\cosh\!\left(\dfrac{L}{2\ell\sqrt{\gamma}}\right) + \sinh\!\left(\dfrac{L}{2\ell\sqrt{\gamma}}\right)}\right],
\end{equation}
which gives $\varepsilon^{\mathrm{V2,unnorm}}(0) \to \frac{1}{\sqrt{\gamma}}\frac{\beta\theta_0}{E}$ as $L/\ell \to \infty$, producing a boundary-layer amplification $\varepsilon^{\mathrm{V2,unnorm}}(0)/\varepsilon^{\mathrm{V1}}(0) \to 1/\sqrt{\gamma}$ at the free ends. Corollary~\ref{cor:artifact} shows that \eqref{eq:unnorm_uniform_strain} is an artifact of kernel truncation at $x=0,L$:

\begin{corollary}[Elimination of the Static $1/\sqrt{\gamma}$ Boundary Layer Under Normalized Kernels]\label{cor:artifact}
Let $\Kcal_\ell^{\mathrm{norm}}$ be any symmetric positive-semidefinite kernel on $[0,L]^2$ satisfying the partition-of-unity condition $\int_0^L \Kcal_\ell^{\mathrm{norm}}(x,\xi)\dd\xi \equiv 1$ (including $W_\ell^{\mathrm{BFP}}$ in \eqref{eq:W_BFP} and $K_\ell^{\mathrm{Neu}}$ in \eqref{eq:K_Neu}). Then under a uniform static temperature change $\theta(x) \equiv \theta_0$ on a traction-free rod $\sigma(x) \equiv 0$, both Variant~V1 and Variant~V2 yield the identical, spatially uniform free thermal expansion everywhere on $[0,L]$:
\begin{equation}\label{eq:normalized_uniform_strain}
  \varepsilon^{\mathrm{V2,norm}}(x) \equiv \varepsilon^{\mathrm{V1,norm}}(x) \equiv \frac{\beta\theta_0}{E} \qquad \forall\, x \in [0,L], \quad \forall\, \gamma \in (0,1).
\end{equation}
\end{corollary}

In particular, substituting the constant field $f(x) = E\varepsilon^{\mathrm{V2}}(x) \equiv \beta\theta_0$ and $g(x) \equiv \beta\theta_0$ into the BFP CBC \eqref{eq:cbc_bfp} gives $g - \gamma_{\mathrm{eff}}f = (1-\gamma)(1-\vartheta_\ell(x))\beta\theta_0$, which satisfies $(\partial_x \mp \ell^{-1})(1-\vartheta_\ell(x))\big|_{x=0,L} = 0$ identically; similarly, under the Neumann CBC \eqref{eq:cbc_neu}, $\partial_x\big((1-\gamma)\beta\theta_0\big)\big|_{x=0,L} = 0$ holds identically. Thus, \textbf{we state plainly that the static $1/\sqrt{\gamma}$ surface thermal-strain layer of \eqref{eq:unnorm_uniform_strain} is a kernel-truncation artifact of $K_\ell^{\mathrm{unnorm}}$, not a physical thermoelastic size effect.} Genuine physical differences between Variant~V1 and Variant~V2 on bounded domains occur only when $\partial_x\theta(x,t) \neq 0$, where $\Lcal_{\gamma,\ell}^{\mathrm{norm}}[\beta\theta] \neq \beta\theta$.

\section{Calibrated Magnitude Estimate in the Fourier-Valid Regime}\label{sec:calibration}

\subsection{Atomistic Calibration of $\ell$ and Non-Identifiability of $\gamma$ in FCC Aluminum}\label{subsec:al_calibration}

To assess the physical magnitude of the bulk dispersion and attenuation differences derived in Section~\ref{sec:mode_count}, we evaluate them for face-centered cubic (fcc) Aluminum (Al) at $T_0 = 300\;\mathrm{K}$ using published, independently calibrated atomistic length scales:
\begin{enumerate}[leftmargin=*,align=left]
  \item \textbf{Atomistic 2NN MEAM calibration (Lazar, Agiasofitou, and B{\"o}hlke \cite{Lazar2022}):} Using the Second Nearest-Neighbor Modified Embedded Atom Method (2NN MEAM) interatomic potential for fcc Aluminum (lattice parameter $a = 4.04950\;\mathrm{\AA} = 0.404950\;\mathrm{nm}$), Lazar, Agiasofitou, and B{\"o}hlke \cite{Lazar2022} computed the sixth-order strain-gradient constitutive tensor and extracted the cubic longitudinal characteristic length $\ell_1 = 1.19303\;\mathrm{\AA} = 0.119303\;\mathrm{nm}$ (Table~3 of \cite{Lazar2022}) and the Voigt-averaged isotropic longitudinal characteristic length (Table~6 of \cite{Lazar2022}):
  \begin{equation}\label{eq:lazar_ell}
    \bar{\ell}_1(\mathrm{Al}) = 1.20272\;\mathrm{\AA} = 0.120272\;\mathrm{nm}, \qquad \frac{\bar{\ell}_1}{a} = 0.2970.
  \end{equation}
  \item \textbf{Brillouin zone-edge lattice dispersion match (Eringen \cite{Eringen1972Linear,Eringen2002}):} Matching the 1D Helmholtz dispersion relation $\omega^2(k) = c_0^2 k^2/(1+\ell^2 k^2)$ to the Born--von K{\'a}rm{\'a}n chain at $k = \pi/a$ gives $e_0 = \sqrt{\pi^2 - 4}/(2\pi) \approx 0.39$, or
  \begin{equation}\label{eq:eringen_ell}
    \ell_{\mathrm{Eringen}}(\mathrm{Al}) = 0.39\,a = 0.157931\;\mathrm{nm}.
  \end{equation}
\end{enumerate}

\paragraph{Degeneracy of $(\gamma,\ell)$ at leading order $O(k^2\ell^2)$.}
Expanding the two-phase symbol \eqref{eq:L_symbol} at long wavelengths ($k\ell \ll 1$) gives $\hat{\Lcal}_{\gamma,\ell}(k^2) = 1 - (1-\gamma)k^2\ell^2 + (1-\gamma)k^4\ell^4 - O(k^6\ell^6)$. Matching the $O(k^2)$ dispersion shift to the first strain-gradient dispersion $\omega^2(k) = c_0^2 k^2(1 - \bar{\ell}_1^2 k^2 + O(k^4))$ fixes only the single effective length
\begin{equation}\label{eq:ell_eff}
  \ell_{\mathrm{eff}} := \sqrt{1-\gamma}\,\ell = \bar{\ell}_1 = 0.120272\;\mathrm{nm}.
\end{equation}
As emphasized by Nobili and Pramanik \cite{NobiliPramanik2025b} (Section~2.1), the two-phase local fraction $\gamma \in (0,1)$ has no independent experimental calibration and acts as a singular-perturbation regularization parameter ensuring coercivity on $L^2(\Omega)$ and boundary-value well-posedness.

\subsection{Scale Separation $\chi_{\mathrm{nl}} = \ell/\ell^*$ and Validity Bounds of Fourier Conduction}\label{subsec:scale_separation}

For fcc Aluminum at $T_0 = 300\;\mathrm{K}$ ($\rho = 2700\;\mathrm{kg/m^3}$, $c_0 = 6420\;\mathrm{m/s}$, $c_E = 900\;\mathrm{J/(kg\cdot K)}$, $\kappa = 237\;\mathrm{W/(m\cdot K)}$, dimensionless thermoelastic coupling $\epsilon = 0.0356$), the characteristic thermoelastic frequency $\omega^*$ and thermal length $\ell^*$ are
\begin{equation}\label{eq:al_thermal_scales}
  \omega^* = \frac{\rho c_E c_0^2}{\kappa} = 4.226 \times 10^{11}\;\mathrm{rad/s} \quad (f^* = \frac{\omega^*}{2\pi} = 67.26\;\mathrm{GHz}), \qquad \ell^* = \frac{c_0}{\omega^*} = 15.19\;\mathrm{nm}.
\end{equation}
Consequently, the dimensionless ratio $\chi_{\mathrm{nl}} := \ell/\ell^*$ between the atomistic nonlocal length and the characteristic thermoelastic length is two orders of magnitude smaller than unity:
\begin{equation}\label{eq:chi_nl_values}
  \chi_{\mathrm{nl}} = \frac{\ell}{\ell^*} = \begin{cases}
    7.917 \times 10^{-3} \quad (\chi_{\mathrm{nl}}^2 = 6.268 \times 10^{-5}), & \text{for } \ell = \bar{\ell}_1 = 0.1203\;\mathrm{nm} \;\text{\cite{Lazar2022}}, \\[4pt]
    1.040 \times 10^{-2} \quad (\chi_{\mathrm{nl}}^2 = 1.081 \times 10^{-4}), & \text{for } \ell = 0.39\,a = 0.1579\;\mathrm{nm} \;\text{\cite{Eringen2002}}.
  \end{cases}
\end{equation}

Furthermore, Fourier's heat conduction law $q = -\kappa\partial_x\theta$ is physically defensible only when (i)~the wave period exceeds the microscopic electron--phonon thermal relaxation time $\tau_q \approx 1.5\;\mathrm{ps}$ ($\omega\tau_q \le 0.1 \implies \Omega := \omega/\omega^* \le 0.158$, $f \le 10.6\;\mathrm{GHz}$) and (ii)~the acoustic wavelength $\lambda = 2\pi/k$ exceeds the thermal carrier mean free path $\Lambda_{\mathrm{th}} \approx 15\;\mathrm{nm}$ by an order of magnitude ($\lambda \ge 10\Lambda_{\mathrm{th}} \implies \Omega \le 0.636$, $f \le 42.8\;\mathrm{GHz}$; cf.\ Guyer and Krumhansl \cite{GuyerKrumhansl1966}). At the strict Fourier validity ceiling $\Omega = 0.158$, we have $k\ell = 1.25 \times 10^{-3}$ and $(k\ell)^2 = 1.56 \times 10^{-6}$; even at the Zener thermoelastic peak $\Omega = 1.0$ ($f = 67.26\;\mathrm{GHz}$, $\lambda = 96.3\;\mathrm{nm}$), we have $k\ell \le 1.03 \times 10^{-2}$ and $(k\ell)^2 \le 1.06 \times 10^{-4}$.

\subsection{Quantitative Comparison Across Frequency Regimes}\label{subsec:quantitative_table}

Table~\ref{tab:al_damping} lists the exact dimensionless quasi-elastic attenuation $\mathrm{Im}(k\ell^*)$ and the relative discrepancies $|\alpha_1^{\mathrm{V1}} - \alpha_1^{\mathrm{V2}}|/\alpha_1^{\mathrm{V1}}$ and $|\alpha_1^{\mathrm{Class\,I}} - \alpha_1^{\mathrm{V1}}|/\alpha_1^{\mathrm{V1}}$ for fcc Aluminum at $300\;\mathrm{K}$ ($\gamma = 0.5$) using the atomistic characteristic length $\bar{\ell}_1 = 0.120272\;\mathrm{nm}$ \cite{Lazar2022}.

\begin{table}[htbp]
  \centering
  \caption{Quasi-elastic wave attenuation comparison for fcc Aluminum at $T_0 = 300\;\mathrm{K}$ ($\epsilon = 0.0356$, $\omega^* = 4.226\times 10^{11}\;\mathrm{rad/s}$, $\ell^* = 15.19\;\mathrm{nm}$, $\bar{\ell}_1 = 0.1203\;\mathrm{nm}$ from Lazar et al.\ \cite{Lazar2022}, $\chi_{\mathrm{nl}} = 7.917\times 10^{-3}$, $\gamma = 0.5$).}\label{tab:al_damping}
  \vspace{4pt}
  \begin{tabular}{ccccccp{4.6cm}}
    \toprule
    $\Omega = \omega/\omega^*$ & $f$ (GHz) & $\lambda$ (nm) & $k\ell$ & $\dfrac{|\alpha_1^{\mathrm{V1}} - \alpha_1^{\mathrm{V2}}|}{\alpha_1^{\mathrm{V1}}}$ & $\dfrac{|\alpha_1^{\mathrm{Class\,I}} - \alpha_1^{\mathrm{V1}}|}{\alpha_1^{\mathrm{V1}}}$ & \textbf{Regime Assessment} \\
    \midrule
    $0.001$ & $0.067$ & $97,137$ & $7.78 \times 10^{-6}$ & $6.50 \times 10^{-11}$ & $1.25 \times 10^{-10}$ & Fourier valid ($\omega\tau_q = 6.3\times 10^{-4}$) \\
    $0.010$ & $0.673$ & $9,714$ & $7.78 \times 10^{-5}$ & $5.32 \times 10^{-9}$ & $1.02 \times 10^{-8}$ & Fourier valid ($\omega\tau_q = 6.3\times 10^{-3}$) \\
    $0.100$ & $6.726$ & $971.2$ & $7.78 \times 10^{-4}$ & $\mathbf{5.34 \times 10^{-7}}$ & $\mathbf{1.02 \times 10^{-6}}$ & \textbf{Strict Fourier ceiling} ($\omega\tau_q = 0.063 \le 0.1$) \\
    $0.500$ & $33.63$ & $193.7$ & $3.90 \times 10^{-3}$ & $\mathbf{1.40 \times 10^{-5}}$ & $\mathbf{2.35 \times 10^{-5}}$ & Pre-ballistic transition ($\omega\tau_q = 0.32$) \\
    $1.000$ & $67.26$ & $96.32$ & $7.85 \times 10^{-3}$ & $\mathbf{5.98 \times 10^{-5}}$ & $\mathbf{7.80 \times 10^{-5}}$ & \textbf{Zener peak} ($\omega\tau_q = 0.63$, diff $< 10^{-4}$) \\
    $10.00$ & $672.6$ & $9.53$ & $7.93 \times 10^{-2}$ & $6.28 \times 10^{-3}$ & $4.79 \times 10^{-3}$ & Fourier invalid ($\omega\tau_q = 6.3 \gg 1$, $\lambda < \Lambda_{\mathrm{th}}$) \\
    $100.0$ & $6,726$ & $0.84$ & $8.98 \times 10^{-1}$ & $6.58 \times 10^{-1}$ & $1.01 \times 10^{0}$ & Continuum \& Fourier invalid ($\lambda \approx 2a$) \\
    \bottomrule
  \end{tabular}
\end{table}

Table~\ref{tab:al_damping} establishes a critical physical conclusion that must be stated plainly: \textbf{within the frequency regime where Fourier heat conduction is physically defensible ($\Omega \lesssim 0.16$--$1.0$), the relative difference in bulk thermoelastic attenuation between Variant~V1, Variant~V2, and Class~I in calibrated fcc Aluminum is of order $O(10^{-6}\text{--}10^{-4})$ ($0.0001\%$ to $0.01\%$), and is identical under Eringen's calibration \eqref{eq:eringen_ell} ($\le 1.35 \times 10^{-4}$ at $\Omega = 1$).} An $O(1)$ divergence between the models occurs only when $k\ell \sim 1$ ($\Omega \sim 100$, $f \sim 6.7\;\mathrm{THz}$, $\lambda \sim 0.84\;\mathrm{nm} \approx 2a$), where the wavelength is comparable to two atomic spacings and both Fourier heat conduction ($\omega\tau_q \approx 63$) and continuum elasticity break down. Accordingly, the importance of enforcing Variant~V1 or Variant~V2 in place of Class~I or Class~II is \emph{mathematical and foundational}---guaranteeing compliance with the Second Law of Thermodynamics, correct spatial differential order ($N=4$ vs.\ $N=3$), and well-posed boundary closure on bounded domains---rather than predicting a large measurable bulk ultrasonic damping split in monolithic crystals.

\section{Limitations and Concluding Remarks}\label{sec:limitations}

\subsection{Limitations of This Note}\label{subsec:limitations}

To avoid any over-interpretation of our results, we record five explicit limitations of the present analysis:
\begin{enumerate}[leftmargin=*,align=left]
  \item \textbf{Convenience sample of the literature:} The census in Section~\ref{subsec:survey} is based on a convenience sample of $N=29$ spatial-nonlocal thermoelasticity papers drawn from a working research collection of 152 mechanics PDFs (of which 13 papers, or $44.8\%$, are co-authored by the present author). It is not a random or exhaustive sample of the entire mechanics literature.
  \item \textbf{One-dimensional isotropic restriction:} To isolate the operator reciprocity, spatial polynomial degree ($N=4$ vs.\ $N=3$), and endpoint Constitutive Boundary Conditions in closed analytical form, the mathematical derivations in Sections~\ref{sec:model}--\ref{sec:boundary} have been carried out for one-dimensional isotropic rods and half-spaces. In three-dimensional anisotropic crystals, the nonlocal kernel is a fourth-order tensor (or a sixth-order strain-gradient tensor \cite{Lazar2022}; see also nonlocal strain-gradient \cite{LimZhangReddy2015} and space-time causal Klein--Gordon nonlocal formulations \cite{LazarAgiasofitou2022KG}) and the boundary closure involves surface-normal and tangential derivatives on curved boundaries $\partial\Omega$.
  \item \textbf{Lack of independent experimental calibration for the two-phase parameter $\gamma$:} As shown in Section~\ref{subsec:al_calibration} and noted by Nobili and Pramanik \cite{NobiliPramanik2025b}, long-wavelength bulk dispersion calibrates only the product $\ell_{\mathrm{eff}}^2 = (1-\gamma)\ell^2 = \bar{\ell}_1^2$. Separating the local phase fraction $\gamma \in (0,1)$ from $\ell$ requires fitting either fourth-order $O(k^4\ell^4)$ lattice dispersion curvature or boundary-layer stiffness in nanoscale resonators, for which no consensus experimental protocol currently exists.
  \item \textbf{Dependence of boundary layers on the choice of normalized kernel:} As demonstrated in Section~\ref{sec:boundary}, while any symmetric partition-of-unity normalized kernel ($\int_\Omega \Kcal_\ell^{\mathrm{norm}}(x,\xi)\dd\xi \equiv 1$, such as $W_\ell^{\mathrm{BFP}}$ or $K_\ell^{\mathrm{Neu}}$) eliminates the unphysical $1/\sqrt{\gamma}$ static thermal-expansion spike of the unnormalized kernel $K_\ell^{\mathrm{unnorm}}$ under uniform heating, the quantitative amplitude of gradient-driven boundary layers under $\partial_x\theta \neq 0$ still depends on whether the boundary attenuates interactions via an additive Dirac delta ($W_\ell^{\mathrm{BFP}}$) or via specular reflection ($K_\ell^{\mathrm{Neu}}$) \cite{BazantLeHoover2010,NobiliPramanik2025b}.
  \item \textbf{Breakdown of Fourier heat conduction before $k\ell \sim 1$ in monolithic crystals:} As shown in Section~\ref{sec:calibration}, in monolithic crystalline metals and semiconductors where $\ell \sim 0.12$--$0.21\;\mathrm{nm}$ and $\ell^* \sim 10$--$15\;\mathrm{nm}$, the ratio $\chi_{\mathrm{nl}} = \ell/\ell^* \sim 10^{-2}$ is small, so nonlocal corrections to bulk thermoelastic damping inside the Fourier-valid regime ($\omega\tau_q \le 0.1$) are $O((k\ell)^2) \lesssim 10^{-4}$. Reaching $k\ell \sim 1$ at sub-GHz or GHz frequencies where diffusive heat transport remains valid would require micro-architectured phononic metamaterials with a mesoscopic unit-cell length $\ell_{\mathrm{meta}} \sim \ell^*$, which lies outside the scope of atomistic continuum elasticity.
\end{enumerate}

\subsection{Summary of Practical Guidelines for Nonlocal Thermoelastic Modeling}\label{subsec:guidelines}

Whenever Eringen's integral nonlocal elasticity is coupled with thermoelasticity and recast into differential form on a domain $\Omega$ with boundaries, three consistency conditions must be satisfied simultaneously:
\begin{enumerate}[leftmargin=*,align=left]
  \item \textbf{Enforce Eringen's \cite{Eringen1974} operator reciprocity between thermal stress and dilational entropy:} Do not pair the stress-nonlocal relation $(1-\ell^2\partial_x^2)\sigma = E\varepsilon - \beta\theta$ with the local constant-strain entropy $\rho\eta = \beta\varepsilon + (\rho c_E/T_0)\theta$ (Theorem~\ref{thm:no_free_energy}). Either apply the two-phase operator $\Lcal_{\gamma,\ell}$ to both the thermal stress and the dilational entropy (\textbf{Variant~V1}, Eqs.~\eqref{eq:V1_diff_stress}--\eqref{eq:V1_diff_heat}, eighth-order in space, $N=4$ mode pairs), or keep both the thermal stress and the dilational entropy local while applying $\Lcal_{\gamma,\ell}$ only to the elastic stiffness (\textbf{Variant~V2}, Eqs.~\eqref{eq:V2_diff_stress}--\eqref{eq:V2_diff_heat}, sixth-order in space, $N=3$ mode pairs). If a single-phase stress-gradient Gibbs model $(1-\ell^2\partial_x^2)\sigma = E\varepsilon - \beta\theta$ is used instead, the entropy must include the Gibbs stress-Laplacian term $\rho\eta^{\mathrm{Gibbs}} = \beta\varepsilon + (\rho c_E/T_0)\theta + (\beta\ell^2/E)\partial_x^2\sigma$ (Eq.~\eqref{eq:gibbs_entropy}).
  \item \textbf{Impose physical traction BCs on the nonlocal Cauchy stress $\sigma$ and close the elevated PDE order with CBCs:} Boundary traction continuity must be imposed on the physical Cauchy stress $\sigma$ that enters the momentum balance $\partial_x\sigma = \rho\ddot{u}$, never on the auxiliary local stress $\sigma^{\mathrm{loc}} = E\varepsilon - \beta\theta$, and the extra spatial order ($N=4$ for V1, $N=3$ for V2) must be closed using the Constitutive Boundary Conditions of Theorem~\ref{thm:cbc}.
  \item \textbf{Use a partition-of-unity normalized symmetric kernel on bounded domains:} To ensure that a traction-free body subjected to a uniform static temperature rise $\theta_0$ undergoes uniform free thermal expansion $\varepsilon(x) \equiv \beta\theta_0/E$ without artificial $1/\sqrt{\gamma}$ surface strain spikes (Corollary~\ref{cor:artifact}), the integral operator on $\Omega$ should be formulated with a symmetric constant-preserving kernel such as $W_\ell^{\mathrm{BFP}}$ \eqref{eq:W_BFP} or $K_\ell^{\mathrm{Neu}}$ \eqref{eq:K_Neu}.
\end{enumerate}


\appendix
\section{Complete Mathematical Proofs}\label{app:proofs}

\subsection{Proof of Proposition \ref{prop:reciprocity} (Master Free Energy and Reciprocity on $L^2(\Omega)$)}\label{app:proof_prop1}

\begin{proof}
Let $\mathcal{H} := L^2(\Omega) \times L^2(\Omega)$ equipped with the inner product $\langle(\varepsilon_1,\theta_1),(\varepsilon_2,\theta_2)\rangle_{\mathcal{H}} = \langle\varepsilon_1,\varepsilon_2\rangle_{L^2(\Omega)} + \langle\theta_1,\theta_2\rangle_{L^2(\Omega)}$. Suppose $\Psi : \mathcal{H} \to \Rbb$ is a twice Fr{\'e}chet-differentiable functional generating the linear state laws \eqref{eq:general_state_laws} via \eqref{eq:frechet_state_laws}:
\begin{align}\label{eq:app_first_variation}
  \mathrm{D}\Psi(\varepsilon,\theta)[(\delta\varepsilon,\delta\theta)]
  &= \langle\sigma,\delta\varepsilon\rangle_{L^2(\Omega)} - \langle\rho\eta,\delta\theta\rangle_{L^2(\Omega)} \nonumber\\
  &= \langle E\Lcal_{\mathrm{el}}\varepsilon - \beta\Lcal_{\mathrm{th}}\theta,\,\delta\varepsilon\rangle_{L^2(\Omega)} - \left\langle\beta\Lcal_\eta\varepsilon + \frac{\rho c_E}{T_0}\Lcal_c\theta,\,\delta\theta\right\rangle_{L^2(\Omega)}.
\end{align}
Taking the second Fr{\'e}chet derivative $\mathrm{D}^2\Psi(\varepsilon,\theta)$ along mixed variations $(\delta\varepsilon,0)$ and $(0,\delta\theta)$ gives
\begin{align}
  \mathrm{D}^2\Psi(\varepsilon,\theta)[(\delta\varepsilon,0),(0,\delta\theta)] &= -\beta\,\langle\Lcal_{\mathrm{th}}\delta\theta,\,\delta\varepsilon\rangle_{L^2(\Omega)} = -\beta\,\langle\delta\theta,\,\Lcal_{\mathrm{th}}^*\delta\varepsilon\rangle_{L^2(\Omega)}, \label{eq:app_d2_1}\\
  \mathrm{D}^2\Psi(\varepsilon,\theta)[(0,\delta\theta),(\delta\varepsilon,0)] &= -\beta\,\langle\Lcal_\eta\delta\varepsilon,\,\delta\theta\rangle_{L^2(\Omega)} = -\beta\,\langle\delta\theta,\,\Lcal_\eta\delta\varepsilon\rangle_{L^2(\Omega)}. \label{eq:app_d2_2}
\end{align}
By Schwarz's theorem on the symmetry of the second Fr{\'e}chet derivative of a $C^2$ functional on a Hilbert space, \eqref{eq:app_d2_1} and \eqref{eq:app_d2_2} must be equal for all $\delta\varepsilon,\delta\theta \in L^2(\Omega)$. Since $\beta > 0$, this holds if and only if $\langle\delta\theta, (\Lcal_\eta - \Lcal_{\mathrm{th}}^*)\delta\varepsilon\rangle_{L^2(\Omega)} = 0$ for all $\delta\varepsilon,\delta\theta \in L^2(\Omega)$, which is equivalent to $\Lcal_\eta = \Lcal_{\mathrm{th}}^*$. Symmetry under pure strain variations and pure thermal variations holds because $\Lcal_{\mathrm{el}}^* = \Lcal_{\mathrm{el}}$ and $\Lcal_c^* = \Lcal_c$.

Conversely, when $\Lcal_{\mathrm{el}} = \Lcal_{\mathrm{el}}^*$, $\Lcal_c = \Lcal_c^*$, and $\Lcal_\eta = \Lcal_{\mathrm{th}}^* = \Lcal_{\mathrm{th}}$, differentiating the quadratic functional $\Psi[\varepsilon,\theta]$ in \eqref{eq:master_psi} along any smooth trajectory $(\varepsilon(\cdot,t),\theta(\cdot,t)) \in C^1([0,T];\mathcal{H})$ and using the self-adjointness identities $\langle\varepsilon,\Lcal_{\mathrm{el}}\dot{\varepsilon}\rangle = \langle\Lcal_{\mathrm{el}}\varepsilon,\dot{\varepsilon}\rangle$, $\langle\theta,\Lcal_c\dot{\theta}\rangle = \langle\Lcal_c\theta,\dot{\theta}\rangle$, and $\langle\varepsilon,\Lcal_{\mathrm{th}}\dot{\theta}\rangle = \langle\Lcal_{\mathrm{th}}\varepsilon,\dot{\theta}\rangle$ yields
\begin{align}
  \frac{\dd}{\dd t}\Psi[\varepsilon,\theta] &= \langle E\Lcal_{\mathrm{el}}\varepsilon - \beta\Lcal_{\mathrm{th}}\theta,\,\dot{\varepsilon}\rangle_{L^2(\Omega)} - \left\langle\beta\Lcal_{\mathrm{th}}\varepsilon + \frac{\rho c_E}{T_0}\Lcal_c\theta,\,\dot{\theta}\right\rangle_{L^2(\Omega)} \nonumber\\
  &= \int_\Omega \left[\sigma(x,t)\dot{\varepsilon}(x,t) - \rho\eta(x,t)\dot{\theta}(x,t)\right]\dd x. \label{eq:app_psi_dot}
\end{align}
Substituting \eqref{eq:app_psi_dot} into the global Clausius--Duhem inequality \eqref{eq:global_cd} cancels the conservative power identically, leaving $\Dcal_{\mathrm{global}}(t) = -\int_\Omega \frac{q}{T_0}\partial_x\theta\dd x = \int_\Omega \frac{\kappa}{T_0}(\partial_x\theta)^2\dd x \ge 0$.
\end{proof}

\subsection{Proof of Theorem \ref{thm:no_free_energy} (Incompatibility of Class I in Helmholtz and Gibbs Representations)}\label{app:proof_thm1}

\begin{proof}
\textbf{Part (a):} For the constitutive pairing \eqref{eq:thm1_stress}--\eqref{eq:thm1_entropy}, the thermal stress operator is $\Lcal_{\mathrm{th}} = \Lcal_{\gamma,\ell}$ ($\gamma \in [0,1)$), whereas the dilational entropy operator in $\rho\eta^{\mathrm{loc}} = \beta\varepsilon + \frac{\rho c_E}{T_0}\theta$ is the identity operator $\Lcal_\eta = \Ical$. By Proposition~\ref{prop:reciprocity}, a twice Fr{\'e}chet-differentiable Helmholtz free energy $\Psi[\varepsilon,\theta]$ exists on $L^2(\Omega) \times L^2(\Omega)$ if and only if $\Lcal_\eta = \Lcal_{\mathrm{th}}^* = \Lcal_{\gamma,\ell}$, i.e.,
\begin{equation}\label{eq:app_defect_op}
  \Ical - \Lcal_{\gamma,\ell} = (1-\gamma)\left(\Ical - \Kcal_\ell\right) = 0 \quad \text{on } L^2(\Omega).
\end{equation}
Since $\gamma < 1$ and $\Kcal_\ell$ is a compact integral operator on $L^2(0,L)$ (and has Fourier symbol $(1+\ell^2 k^2)^{-1} \neq 1$ for $k \neq 0$ on $\Rbb$), $\Ical - \Kcal_\ell \neq 0$. Specifically, for any non-constant eigenmode $\phi_n(x) = \cos(n\pi x/L)$ ($n \ge 1$) of the Neumann Laplacian on $[0,L]$, $((1-\gamma)(\Ical - \Kcal_\ell^{\mathrm{Neu}})\phi_n)(x) = \frac{(1-\gamma)k_n^2\ell^2}{1+k_n^2\ell^2}\phi_n(x) \neq 0$. Thus no Helmholtz free energy functional exists.

\medskip\noindent\textbf{Part (b):} In the complementary $(\sigma,\theta)$ representation on $\Omega = [0,L]$, suppose there exists a functional $G[\sigma,\theta] = \int_0^L g(\sigma,\partial_x\sigma,\theta)\dd x$ generating $\varepsilon = \delta G/\delta\sigma = \frac{1}{E}(\sigma - \ell^2\partial_x^2\sigma) + \frac{\beta}{E}\theta$ (which is equivalent to the $\gamma=0$ Class~I stress law $(1-\ell^2\partial_x^2)\sigma = E\varepsilon - \beta\theta$) and $\rho\eta = \delta G/\delta\theta$. The Euler--Lagrange equations of $G[\sigma,\theta]$ require
\begin{align}\label{eq:app_gibbs_el}
  \frac{\partial g}{\partial\sigma} - \partial_x\!\left(\frac{\partial g}{\partial(\partial_x\sigma)}\right) &= \frac{1}{E}\sigma - \frac{\ell^2}{E}\partial_x^2\sigma + \frac{\beta}{E}\theta, \qquad
  \frac{\partial g}{\partial\theta} - \partial_x\!\left(\frac{\partial g}{\partial(\partial_x\theta)}\right) = \rho\eta.
\end{align}
Integrating the first equation of \eqref{eq:app_gibbs_el} with respect to $(\sigma,\partial_x\sigma)$ determines $g$ up to a purely thermal potential $h(\theta)$:
\begin{equation}\label{eq:app_gibbs_g}
  g(\sigma,\partial_x\sigma,\theta) = \frac{1}{2E}\sigma^2 + \frac{\ell^2}{2E}(\partial_x\sigma)^2 + \frac{\beta}{E}\sigma\theta + h(\theta).
\end{equation}
Matching the local limit $\ell \to 0$ ($\sigma = E\varepsilon - \beta\theta$, $\rho\eta = \beta\varepsilon + \frac{\rho c_E}{T_0}\theta$) fixes $h'(\theta) = \left(\frac{\rho c_E}{T_0} + \frac{\beta^2}{E}\right)\theta = \frac{\rho c_\sigma}{T_0}\theta$. Differentiating \eqref{eq:app_gibbs_g} with respect to $\theta$ and substituting $\sigma = E\varepsilon - \beta\theta + \ell^2\partial_x^2\sigma$ yields
\begin{equation}\label{eq:app_gibbs_eta}
  \rho\eta^{\mathrm{Gibbs}} = \frac{\beta}{E}\sigma + \frac{\rho c_\sigma}{T_0}\theta = \beta\varepsilon + \frac{\rho c_E}{T_0}\theta + \frac{\beta\ell^2}{E}\partial_x^2\sigma = \rho\eta^{\mathrm{loc}} + \frac{\beta\ell^2}{E}\partial_x^2\sigma,
\end{equation}
proving \eqref{eq:gibbs_entropy}.

\medskip\noindent\textbf{Part (c):} The explicit evaluation of the closed-cycle integral \eqref{eq:hc_cycle_work}--\eqref{eq:hc_cycle_total} in Section~\ref{subsec:hand_check} for the smooth periodic fields $\varepsilon(x,t) = \varepsilon_0\cos(kx)\cos(\omega t)$ and $\theta(x,t) = \theta_0\cos(kx)\sin(\omega t)$ ($k = n\pi/L$, $n \ge 1$) holds verbatim under Fourier, Cattaneo, Moore--Gibson--Thompson, or phase-lag heat conduction (since $\partial_x\theta$ depends only on $\theta_0$, whereas the negative mechanical--thermal cross-work term in \eqref{eq:hc_cycle_total} is linear in $\varepsilon_0$ and can be made arbitrarily negative by choosing $\varepsilon_0$ large for any fixed $\theta_0 > 0$).
\end{proof}

\subsection{Proof of Theorem \ref{thm:mode_count} (Spatial Mode Counts and $\gamma \to 0^+$ Singular Limit)}\label{app:proof_thm2}

\begin{proof}
\textbf{Part (a):} Expanding \eqref{eq:P_V1} for Variant~V1 in powers of $q = k^2$ gives $P_{\mathrm{V1}}(q,\omega;\gamma) = \sum_{m=0}^4 a_m^{\mathrm{V1}}(\omega;\gamma)q^m$, where the coefficients of $q^4$ and $q^3$ are
\begin{align}
  a_4^{\mathrm{V1}}(\omega;\gamma) &= \gamma\ell^4 E\kappa, \label{eq:app_a4_V1}\\
  a_3^{\mathrm{V1}}(\omega;\gamma) &= \ell^2 E\kappa(1+\gamma) - \rho\omega^2\ell^4\kappa - \mathrm{i}\omega\rho c_E\gamma\ell^4 E - \mathrm{i}\omega T_0\beta^2\gamma^2\ell^4. \label{eq:app_a3_V1}
\end{align}
Since $\gamma, \ell, E, \kappa > 0$, $a_4^{\mathrm{V1}}(\omega;\gamma) > 0$ for all $\omega > 0$ and $\gamma \in (0,1)$, so $\deg_q P_{\mathrm{V1}} = 4$. Moreover, at $q = -\ell^{-2}$ (the pole of $\hat{\Lcal}_{\gamma,\ell}(q)$ cleared in \eqref{eq:P_V1}), evaluating \eqref{eq:P_V1} gives
\begin{equation}\label{eq:app_no_cancellation}
  P_{\mathrm{V1}}(-\ell^{-2},\omega;\gamma) = -\mathrm{i}\omega T_0\beta^2(-\ell^{-2})(1-\gamma)^2 = \mathrm{i}\omega T_0\beta^2\ell^{-2}(1-\gamma)^2 \neq 0,
\end{equation}
which proves that $(1+\ell^2 q)$ does not divide $P_{\mathrm{V1}}(q,\omega;\gamma)$, so all four roots $q_1,q_2,q_3,q_4$ are genuine roots of the rational dispersion equation \eqref{eq:master_dispersion}. For Variant~V2, expanding \eqref{eq:P_V2} gives $\deg_q P_{\mathrm{V2}} = 3$ with leading coefficient $a_3^{\mathrm{V2}} = \gamma\ell^2 E\kappa > 0$ and $P_{\mathrm{V2}}(-\ell^{-2},\omega;\gamma) = -E\ell^{-2}(1-\gamma)(-\kappa\ell^{-2} - \mathrm{i}\omega\rho c_E) \neq 0$. For Class~I (and the bulk PDEs of Class~II), clearing $(1+\ell^2 q)$ for $\gamma \in (0,1)$ gives $\deg_q P_{\mathrm{Class\,I}}(\cdot,\omega;\gamma) = 3$ with $a_3^{\mathrm{Class\,I}} = \gamma\ell^2 E\kappa > 0$; setting $\gamma = 0$ annihilates the $q^3$ term and yields the quadratic \eqref{eq:P_ClassI} ($\deg_q = 2$) with leading coefficient $a_2^{\mathrm{Class\,I}}(\omega;0) = \kappa(E - \rho\omega^2\ell^2) \neq 0$ for $\omega \neq c_0/\ell$.

\medskip\noindent\textbf{Part (b):} As $\gamma \to 0^+$, the coefficient $a_4^{\mathrm{V1}} = \gamma\ell^4 E\kappa$ of $q^4$ in \eqref{eq:P_V1} tends to zero, while $a_3^{\mathrm{V1}}(\omega;0) = \ell^2\kappa(E - \rho\omega^2\ell^2) \neq 0$ (for $\omega \neq c_0/\ell$). By standard singular perturbation of algebraic equations, introducing the rescaled variable $Q := \gamma\ell^2 q$ into $P_{\mathrm{V1}}(Q/(\gamma\ell^2),\omega;\gamma) = 0$ and multiplying by $\gamma^3\ell^6$ yields
\begin{align}\label{eq:app_rescaled_Q}
  &\left[E Q(1+Q) - \gamma\rho\omega^2\ell^2(\gamma+Q)\right]\left(\kappa Q - \mathrm{i}\gamma\omega\rho c_E\ell^2\right)(\gamma+Q) \nonumber\\
  &\qquad - \mathrm{i}\gamma^2\omega T_0\beta^2\ell^2 Q(1+Q)^2 = 0.
\end{align}
At $\gamma = 0$, \eqref{eq:app_rescaled_Q} reduces to $E\kappa Q^3(1+Q) = 0$, which has the triple root $Q = 0$ (corresponding to the three finite roots $q_1,q_2,q_3$ of the $\gamma = 0$ cubic \eqref{eq:P_V1_gamma0}) and the simple root $Q_4(0) = -1$. Applying the implicit function theorem at $Q_4(0) = -1$ gives the Taylor expansion $Q_4(\gamma) = -1 + \frac{\rho\omega^2\ell^2}{E}\gamma - \frac{\mathrm{i}\omega T_0\beta^2\ell^2}{\kappa E}\gamma^2 + O(\gamma^3)$, which divides by $\gamma\ell^2$ to yield \eqref{eq:q4_asymptotic}.

\medskip\noindent\textbf{Part (c):} Writing \eqref{eq:P_V1} and \eqref{eq:P_V2} in terms of the dimensionless variables $\Omega = \omega/\omega^*$, $K = k\ell^*$, $\bar{Q} = K^2 = q(\ell^*)^2$, and $\chi_{\mathrm{nl}} = \ell/\ell^*$, and expanding the quasi-elastic root $\bar{Q}_1(\Omega) = \frac{\Omega^2}{\gamma}\big[1 + \frac{c_1}{\Omega} + \frac{c_2}{\Omega^2} + O(\Omega^{-3})\big]$ as $\Omega \to \infty$ (where $\hat{\Lcal}_{\gamma,\ell}(\bar{Q}_1) = \gamma + O(\Omega^{-2})$) yields $K_1^{\mathrm{V1}}(\Omega) = \frac{\Omega}{\sqrt{\gamma}} + \mathrm{i}\frac{\epsilon}{2}\gamma^{3/2} + O(\Omega^{-1})$ for Variant~V1 and $K_1^{\mathrm{V2}}(\Omega) = \frac{\Omega}{\sqrt{\gamma}} + \mathrm{i}\frac{\epsilon}{2}\gamma^{-1/2} + O(\Omega^{-1})$ for Variant~V2, establishing \eqref{eq:high_freq_plateau}. Expanding for $\chi_{\mathrm{nl}}\Omega \ll 1$ and $\epsilon \ll 1$ similarly gives \eqref{eq:low_k_damping_expansion}.
\end{proof}

\subsection{Proof of Theorem \ref{thm:cbc} and Corollary \ref{cor:artifact} (Boundary Closure and Truncation Artifact)}\label{app:proof_thm3}

\begin{proof}
Let $f \in H^2(0,L)$ and define the auxiliary nonlocal convolution $w(x) := \int_0^L \Kcal_\ell^{\mathrm{base}}(x,\xi)f(\xi)\dd\xi$, where $\Kcal_\ell^{\mathrm{base}} = K_\ell^{\mathrm{unnorm}}$ for Kernels~(i) and~(ii), and $\Kcal_\ell^{\mathrm{base}} = K_\ell^{\mathrm{Neu}}$ for Kernel~(iii).
\begin{enumerate}[label=(\alph*),leftmargin=*]
  \item For $\Kcal_\ell^{\mathrm{base}} = K_\ell^{\mathrm{unnorm}}(x,\xi) = \frac{1}{2\ell}\mathrm{e}^{-|x-\xi|/\ell}$, splitting the integral at $\xi = x$ gives
  \begin{equation}\label{eq:app_w_split}
    w(x) = \frac{1}{2\ell}\mathrm{e}^{-x/\ell}\int_0^x \mathrm{e}^{\xi/\ell}f(\xi)\dd\xi + \frac{1}{2\ell}\mathrm{e}^{x/\ell}\int_x^L \mathrm{e}^{-\xi/\ell}f(\xi)\dd\xi.
  \end{equation}
  Differentiating \eqref{eq:app_w_split} once and twice with respect to $x \in (0,L)$ yields
  \begin{equation}\label{eq:app_w_ode}
    w'(x) = -\frac{1}{2\ell^2}\mathrm{e}^{-x/\ell}\int_0^x \mathrm{e}^{\xi/\ell}f(\xi)\dd\xi + \frac{1}{2\ell^2}\mathrm{e}^{x/\ell}\int_x^L \mathrm{e}^{-\xi/\ell}f(\xi)\dd\xi,
  \end{equation}
  and $w(x) - \ell^2 w''(x) = f(x)$. Evaluating $w(x)$ and $w'(x)$ at $x = 0$ and $x = L$ eliminates the remaining integrals and gives the exact endpoint Robin conditions $w'(0) - \ell^{-1}w(0) = 0$ and $w'(L) + \ell^{-1}w(L) = 0$ \cite{FernandezSaez2016,Romano2017}.
  For Kernel~(i), $g = \gamma f + (1-\gamma)w \iff w = \frac{g - \gamma f}{1-\gamma}$; substituting $w$ into \eqref{eq:app_w_ode} and the Robin conditions yields \eqref{eq:ode_unnorm}--\eqref{eq:cbc_unnorm}.
  \item For Kernel~(ii) ($W_\ell^{\mathrm{BFP}}$), \eqref{eq:L_BFP} gives $g = \gamma_{\mathrm{eff}}(x)f + (1-\gamma)w \iff w = \frac{g - \gamma_{\mathrm{eff}}(x)f}{1-\gamma}$ with the same $w = K_\ell^{\mathrm{unnorm}} * f$. Substituting $w = \frac{g - \gamma_{\mathrm{eff}}f}{1-\gamma}$ into $(1-\ell^2\partial_x^2)w = f$ and $(\partial_x \mp \ell^{-1})w\big|_{0,L} = 0$ yields \eqref{eq:ode_bfp}--\eqref{eq:cbc_bfp}.
  \item For Kernel~(iii) ($K_\ell^{\mathrm{Neu}}$), splitting the integral of \eqref{eq:K_Neu} at $\xi = x$ and differentiating twice gives $(1-\ell^2\partial_x^2)w = f$ on $(0,L)$, while $\partial_x K_\ell^{\mathrm{Neu}}(0,\xi) = \partial_x K_\ell^{\mathrm{Neu}}(L,\xi) = 0$ implies $w'(0) = w'(L) = 0$. Substituting $w = \frac{g - \gamma f}{1-\gamma}$ yields \eqref{eq:ode_neu}--\eqref{eq:cbc_neu}.
\end{enumerate}
Conversely, because $(1-\ell^2\partial_x^2)$ is strictly positive-definite under both the Robin boundary conditions $(\partial_x \mp \ell^{-1})w\big|_{0,L} = 0$ ($\langle(1-\ell^2\partial_x^2)w,w\rangle_{L^2} = \|w\|_{L^2}^2 + \ell^2\|w'\|_{L^2}^2 + \ell[w(0)^2 + w(L)^2] > 0$) and the Neumann boundary conditions $w'(0) = w'(L) = 0$ ($\langle(1-\ell^2\partial_x^2)w,w\rangle_{L^2} = \|w\|_{L^2}^2 + \ell^2\|w'\|_{L^2}^2 > 0$), the homogeneous boundary-value problem has only the trivial solution $w \equiv 0$, so the ODE + CBC system uniquely recovers the integral convolution.

Finally, for Corollary~\ref{cor:artifact}, when $\theta(x) \equiv \theta_0$ is constant and $\Kcal_\ell^{\mathrm{norm}}$ satisfies $\int_\Omega \Kcal_\ell^{\mathrm{norm}}(x,\xi)\dd\xi \equiv 1$, the constant strain field $\varepsilon^{\mathrm{V2}}(x) \equiv \frac{\beta\theta_0}{E}$ satisfies $\Lcal_{\gamma,\ell}^{\mathrm{norm}}[E\varepsilon^{\mathrm{V2}}](x) = \beta\theta_0\big[\gamma + (1-\gamma)\int_\Omega \Kcal_\ell^{\mathrm{norm}}(x,\xi)\dd\xi\big] \equiv \beta\theta_0$. By the $L^2(\Omega)$ coercivity \eqref{eq:self_adjoint_coercive} ($\langle\Lcal_{\gamma,\ell}^{\mathrm{norm}}h,h\rangle \ge \gamma\|h\|_{L^2}^2$), the solution of $\Lcal_{\gamma,\ell}^{\mathrm{norm}}[E\varepsilon^{\mathrm{V2}}] = \beta\theta_0$ in $L^2(\Omega)$ is unique, proving \eqref{eq:normalized_uniform_strain}.
\end{proof}


\bibliographystyle{plain}
\bibliography{references}

\end{document}
